\section{\texorpdfstring{$L$}{L}-symbols of invariant states and permuted blocks}

\subsection{\texorpdfstring{$L$}{L}-symbols and their scalar gauges}

\begin{definition}[L-symbols and compatibility]
    \label{def:mpug_l_symbols}
    \lean{TNLean.Algebra.LSymbol,
        TNLean.Algebra.LSymbol.IsCompatible}
    \leanok
    \uses{def:mpug_scalar_cochains}
    Let $G$ act on a set $\mathsf X$. An $L$-symbol is a function
    $L\colon \mathsf X\times G^2\to\C^\times$, written $L^x_{g,h}$.
    It is compatible with a scalar $3$-cochain $\omega$ when, for every
    $x\in\mathsf X$ and $g,h,k\in G$,
    \begin{align}
        L^x_{g,hk}L^x_{h,k}
         & =\omega(g,h,k)L^{k x}_{g,h}L^x_{gh,k}.
        \label{eq:mpug_l_compatibility}
    \end{align}
    The $L$-symbols were introduced by
    \cite{GarreRubio2023MPOSym}; this is the scalar compatibility relation of
    \cite{FrancoRubio2025MPUGauging}. No transitivity or finiteness assumption
    is imposed on the action.
    % Source anchor: arXiv:2502.20257, eq:omega_and_Ls.
\end{definition}

\begin{definition}[Joint scalar gauge of L-symbols]
    \label{def:mpug_l_gauge}
    \lean{TNLean.Algebra.ActionTensorGauge,
        TNLean.Algebra.LSymbol.gauge}
    \leanok
    \uses{def:mpug_l_symbols, def:mpug_fusion_gauge}
    An action-tensor scalar gauge is a function
    $\gamma\colon G\times\mathsf X\to\C^\times$, written $\gamma^x_g$.
    Together with a fusion-tensor scalar gauge $\beta\colon G^2\to\C^\times$,
    it changes an $L$-symbol by
    \begin{align}
        (\beta,\gamma)\mathbin{\triangleright}L^x_{g,h}
         & =\frac{\gamma^x_{gh}\beta_{g,h}}
        {\gamma^{h x}_g\gamma^x_h}L^x_{g,h}.
        \label{eq:mpug_l_gauge}
    \end{align}
    This is the joint fusion-tensor and action-tensor gauge freedom in
    \cite{FrancoRubio2025MPUGauging}, written with the block label as a
    superscript throughout, as that source does for its block-dependent
    cochains.
    % Source anchor: arXiv:2502.20257, eq:Lsymbgauge (which spells the block
    %   label as a second subscript); the superscript spelling gamma^x_g is
    %   the paper's own at lines 5908 and 5988--5990.
\end{definition}

\begin{lemma}[Identity and composition of joint scalar gauges]
    \label{thm:mpug_l_gauge_composition}
    \lean{TNLean.Algebra.LSymbol.gauge_one,
        TNLean.Algebra.LSymbol.gauge_comp}
    \leanok
    \uses{def:mpug_l_gauge}
    The constant fusion and action gauges leave every $L$-symbol unchanged.
    Successive joint gauges compose by pointwise multiplication:
    \begin{align}
        (\delta,\varepsilon)\mathbin{\triangleright}
        \bigl((\beta,\gamma)\mathbin{\triangleright}L\bigr)
         & = (\beta\delta,\gamma\varepsilon)
        \mathbin{\triangleright}L.
        \label{eq:mpug_l_gauge_composition}
    \end{align}
\end{lemma}

\begin{proof}\leanok
    \uses{def:mpug_l_gauge}
    Substituting (\ref{eq:mpug_l_gauge}) twice and cancelling the intermediate
    factors yields the composition law, and the identity case
    follows by taking all four cochains equal to $1$.
\end{proof}

\begin{definition}[Normalized L-symbols and action-tensor gauges]
    \label{def:mpug_l_normalized}
    \lean{TNLean.Algebra.LSymbol.IsNormalized,
        TNLean.Algebra.ActionTensorGauge.IsNormalized}
    \leanok
    \uses{def:mpug_l_symbols, def:mpug_l_gauge}
    As in \cite{FrancoRubio2025MPUGauging}, the $L$-symbol is normalized when,
    for every $x\in\mathsf X$ and $g\in G$,
    \begin{align}
        L^x_{g,e} & =1.
        \label{eq:mpug_l_normalized_right}
    \end{align}
    \begin{align}
        L^x_{e,g} & =1.
        \label{eq:mpug_l_normalized_left}
    \end{align}
    Alongside (\ref{eq:mpug_l_normalized_right}) and
    (\ref{eq:mpug_l_normalized_left}), the action-tensor gauge is normalized
    when
    \begin{align}
        \gamma^x_e & =1
        \label{eq:mpug_action_gauge_normalized}
    \end{align}
    for every $x\in\mathsf X$. This last condition is not a displayed
    equation of the source: it is the scalar shadow of the standing
    convention that every action tensor involving the identity is trivial.
    % Source anchor: arXiv:2502.20257, eq:triv_Ls for the L-symbol axes;
    %   gamma^x_e = 1 is not a labelled source equation but the standing
    %   convention at lines 1936--1937 that any fusion or action tensor
    %   involving e is trivial, applied to the action-tensor gauge scalar of
    %   eq:scalar_act_ten.
\end{definition}

\begin{definition}[Restricted scalar gauges]
    \label{def:mpug_restricted_scalar_gauge}
    \lean{TNLean.Algebra.ScalarCocycle.IsInverseNormalized,
        TNLean.Algebra.RestrictedFusionGauge,
        TNLean.Algebra.RestrictedScalarGauge}
    \leanok
    \discussion{7700}
    \uses{def:mpug_fusion_gauge, def:mpug_l_gauge, def:mpug_l_normalized}
    A fusion scalar $\beta\colon G^2\to\C^\times$ is
    \emph{inverse-normalized} when, for every $g\in G$,
    \begin{align}
        \beta_{g,e}=\beta_{e,g}=\beta_{g,g^{-1}}
         & =1.
        \label{eq:mpug_restricted_fusion_gauge}
    \end{align}
    Following the standing restriction imposed at
    \cite[lines~2050--2054]{FrancoRubio2025MPUGauging} after fixing the
    fusion and action gauges of Definition~\ref{def:mpug_l_gauge}, a
    restricted scalar gauge is a pair $(\beta,\gamma)$ in which $\beta$ is
    inverse-normalized and, for every $x\in\mathsf X$,
    \begin{align}
        \gamma^x_e & =1.
        \label{eq:mpug_restricted_action_gauge}
    \end{align}
    % Source anchor: arXiv:2502.20257, unlabelled display at lines 2050--2054,
    % immediately after eq:zetas_in_gauge.
\end{definition}

\begin{lemma}[Joint scalar gauges preserve L-symbol compatibility]
    \label{thm:mpug_l_gauge_preservation}
    \lean{TNLean.Algebra.LSymbol.IsCompatible.gauge,
        TNLean.Algebra.LSymbol.gauge_apply_right_one,
        TNLean.Algebra.LSymbol.gauge_apply_left_one,
        TNLean.Algebra.LSymbol.isNormalized_gauge_iff,
        TNLean.Algebra.LSymbol.IsNormalized.gauge}
    \leanok
    \uses{def:mpug_fusion_gauge, def:mpug_l_symbols,
        def:mpug_l_gauge, def:mpug_l_normalized}
    If $L$ is compatible with $\omega$, then
    $(\beta,\gamma)\mathbin{\triangleright}L$ is compatible with
    $\beta\mathbin{\triangleright}\omega$. On the identity axes, the gauge
    formula is
    \begin{align}
        ((\beta,\gamma)\mathbin{\triangleright}L)^x_{g,e}
         & =\frac{\beta_{g,e}}{\gamma^x_e}L^x_{g,e}.
        \label{eq:mpug_l_gauge_right}
    \end{align}
    \begin{align}
        ((\beta,\gamma)\mathbin{\triangleright}L)^x_{e,g}
         & =\frac{\beta_{e,g}}{\gamma^{g x}_e}L^x_{e,g}.
        \label{eq:mpug_l_gauge_left}
    \end{align}
    Thus the gauged $L$-symbol is normalized exactly when
    $\beta_{g,e}L^x_{g,e}=\gamma^x_e$ and
    $\beta_{e,g}L^x_{e,g}=\gamma^{g x}_e$ for all $x$ and $g$. In particular,
    normalized $\beta$, $\gamma$, and $L$ give a normalized gauged $L$-symbol.
\end{lemma}

\begin{proof}\leanok
    \uses{def:mpug_fusion_gauge, def:mpug_l_symbols, def:mpug_l_gauge,
        def:mpug_l_normalized}
    With $L'=(\beta,\gamma)\mathbin{\triangleright}L$, substituting
    (\ref{eq:mpug_l_gauge}) into the left side of
    (\ref{eq:mpug_l_compatibility}) for $L'$ cancels the factor
    $\gamma^x_{hk}$ and leaves
    \begin{align}
        {L'}^x_{g,hk}{L'}^x_{h,k}
         & =\frac{\gamma^x_{ghk}}
        {\gamma^{(hk)x}_g\gamma^{k x}_h\gamma^x_k}\,
        \beta_{g,hk}\beta_{h,k}\,L^x_{g,hk}L^x_{h,k}.
        \label{eq:mpug_l_gauge_left_side}
    \end{align}
    The same substitution on the right side cancels
    $\gamma^{k x}_{gh}$ and leaves
    \begin{align}
        {L'}^{k x}_{g,h}{L'}^x_{gh,k}
         & =\frac{\gamma^x_{ghk}}
        {\gamma^{h(k x)}_g\gamma^{k x}_h\gamma^x_k}\,
        \beta_{g,h}\beta_{gh,k}\,L^{k x}_{g,h}L^x_{gh,k}.
        \label{eq:mpug_l_gauge_right_side}
    \end{align}
    The action law $h(k x)=(hk)x$ makes the two action-tensor prefactors of
    (\ref{eq:mpug_l_gauge_left_side}) and
    (\ref{eq:mpug_l_gauge_right_side}) equal. Applying
    (\ref{eq:mpug_l_compatibility}) for $L$ and $\omega$ to the right of
    (\ref{eq:mpug_l_gauge_left_side}) therefore gives
    \begin{align}
        {L'}^x_{g,hk}{L'}^x_{h,k}
         & =\frac{\beta_{g,hk}\beta_{h,k}}
        {\beta_{g,h}\beta_{gh,k}}\,\omega(g,h,k)\,
        {L'}^{k x}_{g,h}{L'}^x_{gh,k},
        \label{eq:mpug_l_gauge_compatibility}
    \end{align}
    whose scalar factor is
    $(\beta\mathbin{\triangleright}\omega)(g,h,k)$ by
    (\ref{eq:mpug_coboundary}) and
    (\ref{eq:mpug_gauge_action}). Setting either group argument equal
    to $e$ in (\ref{eq:mpug_l_gauge}) gives
    (\ref{eq:mpug_l_gauge_right}) and
    (\ref{eq:mpug_l_gauge_left}), from which the normalization statements
    follow.
\end{proof}

\begin{lemma}[A block fixed by a subgroup trivializes the restricted three-cocycle]
    \label{thm:mpug_fixed_block_trivializes}
    \lean{TNLean.Algebra.LSymbol.isTrivialGaugeClass_comap_of_isCompatible,
        TNLean.Algebra.LSymbol.isTrivialGaugeClass_of_isCompatible_of_forall_smul_eq}
    \leanok
    \uses{def:mpug_l_symbols, def:mpug_cochain_restriction,
        def:mpug_fusion_gauge, def:mpug_trivial_gauge_class}
    Let $L$ be compatible with $\omega$, let $f\colon H\to G$ be a group
    homomorphism, and let $x\in\mathsf X$ satisfy $f(a)x=x$ for every
    $a\in H$. Then $f^*\omega$ has trivial gauge class; explicitly,
    $f^*\omega=d\lambda$ with $\lambda(a,b)=L^x_{f(a),f(b)}$. In particular,
    if some $x$ is fixed by all of $G$, then $\omega$ has trivial gauge class.
    This provides the scalar step both for the statement of
    \cite[lines~740--745]{GarreRubio2023MPOSym} that a subgroup leaving a block
    invariant trivializes the three-cocycle and for the statement of
    \cite[line~1925]{FrancoRubio2025MPUGauging} that a single invariant block
    reduces (\ref{eq:mpug_l_compatibility}) to the three-coboundary condition.
\end{lemma}

\begin{proof}\leanok
    \uses{def:mpug_l_symbols, def:mpug_fusion_gauge,
        def:mpug_cochain_restriction, def:mpug_trivial_gauge_class}
    For $a,b,c\in H$, with $g=f(a)$, $h=f(b)$, and $k=f(c)$, the invariance
    $kx=x$ turns (\ref{eq:mpug_l_compatibility}) into
    \begin{align}
        \omega(g,h,k)
         & =\frac{L^x_{g,hk}L^x_{h,k}}{L^x_{g,h}L^x_{gh,k}},
    \end{align}
    which is $(d\lambda)(a,b,c)$ by (\ref{eq:mpug_coboundary}) and the
    multiplicativity of $f$. Thus
    $f^*\omega=\lambda\mathbin{\triangleright}1$. For a block fixed by all
    of $G$, taking $f$ to be the identity completes the proof.
\end{proof}

\begin{lemma}[Action of the restricted scalar gauge group]
    \label{thm:mpug_restricted_scalar_gauge_action}
    \lean{TNLean.Algebra.ScalarCocycle.isInverseNormalized_one,
        TNLean.Algebra.ScalarCocycle.IsInverseNormalized.mul,
        TNLean.Algebra.ScalarCocycle.IsInverseNormalized.inv,
        TNLean.Algebra.ActionTensorGauge.isNormalized_one,
        TNLean.Algebra.ActionTensorGauge.IsNormalized.mul,
        TNLean.Algebra.ActionTensorGauge.IsNormalized.inv,
        TNLean.Algebra.RestrictedFusionGauge.one,
        TNLean.Algebra.RestrictedFusionGauge.mul,
        TNLean.Algebra.RestrictedFusionGauge.inv,
        TNLean.Algebra.RestrictedScalarGauge.one,
        TNLean.Algebra.RestrictedScalarGauge.mul,
        TNLean.Algebra.RestrictedScalarGauge.inv,
        TNLean.Algebra.RestrictedFusionGauge.fusionGauge_one,
        TNLean.Algebra.RestrictedFusionGauge.isNormalized_fusionGauge,
        TNLean.Algebra.RestrictedFusionGauge.isCocycle_fusionGauge,
        TNLean.Algebra.RestrictedFusionGauge.fusionGauge_mul,
        TNLean.Algebra.RestrictedFusionGauge.fusionGauge_inv,
        TNLean.Algebra.RestrictedScalarGauge.gauge_one,
        TNLean.Algebra.RestrictedScalarGauge.isNormalized_gauge,
        TNLean.Algebra.RestrictedScalarGauge.isCompatible_gauge,
        TNLean.Algebra.RestrictedScalarGauge.gauge_mul,
        TNLean.Algebra.RestrictedScalarGauge.gauge_inv}
    \leanok
    \discussion{7700}
    \uses{def:mpug_restricted_scalar_gauge}
    Let $\kappa=(\beta,\gamma)$ be a restricted scalar gauge. If $\omega$ is
    a normalized scalar $3$-cocycle and $L$ is a normalized $L$-symbol
    compatible with $\omega$, then
    \begin{align}
        \omega' & =\beta\mathbin{\triangleright}\omega,     \\
        L'      & =(\beta,\gamma)\mathbin{\triangleright}L.
        \label{eq:mpug_restricted_gauge_action}
    \end{align}
    are again normalized, $\omega'$ is a scalar $3$-cocycle, and $L'$ is
    compatible with $\omega'$. The identity gauge acts trivially, successive
    restricted gauges compose by pointwise multiplication, and applying
    $\kappa^{-1}$ after $\kappa$ recovers $(\omega,L)$.
    These are scalar statements only.
    % The theorem does not choose fusion tensors, action tensors, or
    % tensor-level gauges satisfying the same normalization. Tensor-level
    % restricted gauge choices remain issue #7325.
\end{lemma}

\begin{proof}\leanok
    \uses{def:mpug_restricted_scalar_gauge,
        thm:mpug_gauge_preserves_cocycle,
        thm:mpug_gauge_preserves_normalization,
        thm:mpug_l_gauge_composition,
        thm:mpug_l_gauge_preservation}
    The three specified values of $\beta$ and the identity value of $\gamma$
    remain equal to $1$ under pointwise products and inverses, so the general
    preservation laws apply to~(\ref{eq:mpug_restricted_gauge_action}).
    The group laws and the composition formulas for the two gauge actions then
    give the identity, composition, and inverse assertions.
\end{proof}

\subsection{An anomaly obstructs invariant normal states}

\begin{definition}[L-symbol of an invariant normal state]
    \label{def:mpug_single_block_l_symbol}
    \lean{MPOTensor.GroupFamily.FixesMPV,
        MPOTensor.GroupFamily.ActionData,
        MPOTensor.GroupFamily.ActionData.actV,
        MPOTensor.GroupFamily.ActionData.fuseV,
        MPOTensor.GroupFamily.ActionData.lSymbol}
    \leanok
    \uses{def:mpug_anomaly_three_cochain, def:mpug_l_symbols, def:mpv_state}
    Let $\{O_g\}$ be as in Definition~\ref{def:mpug_anomaly_three_cochain},
    and let $A$ be a matrix product state tensor with
    $O_g\ket{V^{(N)}(A)}=\ket{V^{(N)}(A)}$ for every $g$ and every $N\geq1$.
    Action tensors for $A$ are, for every $g$, a reduction $(V_g,W_g)$ of
    $O_g\cdot A$ onto $A$. The stacked action $(O_gO_h)\cdot A$ then has the
    two left boundaries
    \begin{align}
        \alpha_{g,h}  & =V_g(\mathbf 1\otimes V_h)\,\mathsf a^{-1}, &
        \varphi_{g,h} & =V_{gh}(F^<_{g,h}\otimes\mathbf 1),
        \notag
    \end{align}
    and the $L$-symbol of a single invariant block in
    \cite{FrancoRubio2025MPUGauging} is a nonzero scalar $L_{g,h}=z$ with
    $\alpha_{g,h}B^{\mathbf w}=z\,\varphi_{g,h}B^{\mathbf w}$ for all long
    words $\mathbf w$ of $B=(O_gO_h)\cdot A$, or $1$ if there is none.
\end{definition}

\begin{center}
    % Formula: \alpha_{g,h}B^{\mathbf w}=L_{g,h}\,\varphi_{g,h}B^{\mathbf w}
    % with \alpha_{g,h}=V_g(\mathbf 1\otimes V_h)\mathsf a^{-1},
    % \varphi_{g,h}=V_{gh}(F^<_{g,h}\otimes\mathbf 1), and
    % B=(O_gO_h)\cdot A; source: arXiv:2502.20257, eq:defL, lines 1875--1914
    % (drawn there with right boundaries).
    % Ink-to-index map: the bracketed column stands for the word \mathbf w
    % of B, with rows O_g, O_h and the state tensor A from top to bottom;
    % each pill has its single side west.  On the left, V_h reduces
    % O_h\cdot A and V_g then reduces O_g\cdot A; on the right, F^<_{g,h}
    % fuses O_gO_h and V_{gh} reduces O_{gh}\cdot A.  The associator is
    % not drawn.
    % Contraction set: the physical indices between the rows, the bonds
    % between the pills, and the bonds into the word; the three east virtual
    % legs, the west leg, and the upper physical leg remain open.
    % Boundary signature: (virtual-west=1 | virtual-east=3 |
    % physical-up=1 | physical-down=0) on both sides.
    \tenkzkernel
    \begin{tenkz}[rows={op,op,ket}, cols=3, bonds=none]
        \tnfuse[at=(1,1), name=Fo, skin=pill, wires=3]{V_g}
        \tnfuse[at=(2,2), name=Fi, skin=pill]{V_h}
        \tn[at=(1,2), void=sealed]{}
        \tn[at=(1,3), name=Tg, skin=mpo, ports={90:physical, 270:physical, 180:virtual, 0:virtual}]{O_g}
        \tn[at=(2,3), name=Th, skin=mpo, ports={90:physical, 270:physical, 180:virtual, 0:virtual}]{O_h}
        \tn[at=(3,3), name=A, label pos=s, ports={90:physical, 180:virtual, 0:virtual}]{A}
        \tnwire{Tg.270}{Th.90}
        \tnwire{Th.270}{A.90}
        \tnwire{open w}{Fo.180}
        \tnwire{Fo.0@1}{Tg.180}
        \tnwire{Fo.0@2}{Fi.180}
        \tnwire{Fi.0@1}{Th.180}
        \tnwire{Fi.0@2}{A.180}
        \tnwire{Tg.0}{open e}
        \tnwire{Th.0}{open e}
        \tnwire{A.0}{open e}
        \tnmark[form=bracket, label pos=90]{Tg}{\mathbf w}
    \end{tenkz}
    \quad $=L_{g,h}$ \quad
    \begin{tenkz}[rows={op,op,ket}, cols=3, bonds=none]
        \tnfuse[at=(1,1), name=Fo, skin=pill, wires=3]{V_{gh}}
        \tnfuse[at=(1,2), name=Fi, skin=pill]{F^<_{g,h}}
        \tn[at=(3,2), void=sealed]{}
        \tn[at=(1,3), name=Tg, skin=mpo, ports={90:physical, 270:physical, 180:virtual, 0:virtual}]{O_g}
        \tn[at=(2,3), name=Th, skin=mpo, ports={90:physical, 270:physical, 180:virtual, 0:virtual}]{O_h}
        \tn[at=(3,3), name=A, label pos=s, ports={90:physical, 180:virtual, 0:virtual}]{A}
        \tnwire{Tg.270}{Th.90}
        \tnwire{Th.270}{A.90}
        \tnwire{open w}{Fo.180}
        \tnwire{Fo.0@1}{Fi.180}
        \tnwire{Fo.0@3}{A.180}
        \tnwire{Fi.0@1}{Tg.180}
        \tnwire{Fi.0@2}{Th.180}
        \tnwire{Tg.0}{open e}
        \tnwire{Th.0}{open e}
        \tnwire{A.0}{open e}
        \tnmark[form=bracket, label pos=90]{Tg}{\mathbf w}
    \end{tenkz}
\end{center}

\begin{lemma}[Characterization of the single-block L-symbol]
    \label{thm:mpug_single_block_l_symbol_spec}
    \lean{MPOTensor.GroupFamily.nonempty_actionData,
        MPOTensor.GroupFamily.ActionData.isDressedProportional_lSymbol}
    \leanok
    \uses{def:mpug_single_block_l_symbol}
    If $A$ in Definition~\ref{def:mpug_single_block_l_symbol} is normal
    of positive bond dimension, then action tensors for $A$ exist
    \cite[lines~1064--1090]{GarreRubio2023MPOSym}, and for every choice of
    them $\alpha_{g,h}B^{\mathbf w}=L_{g,h}\,\varphi_{g,h}B^{\mathbf w}$ for
    all long words $\mathbf w$ of $B=(O_gO_h)\cdot A$.
\end{lemma}

\begin{proof}\leanok
    \uses{cor:asym_action_tensors, thm:mpug_action_tensor_reductions,
        thm:mpug_reduction_composition, thm:mpug_dressed_proportionality}
    Existence is Corollary~\ref{cor:asym_action_tensors}. By
    Lemmas~\ref{thm:mpug_action_tensor_reductions}
    and~\ref{thm:mpug_reduction_composition}, both $\alpha_{g,h}$ and
    $\varphi_{g,h}$ are left boundaries of reductions of $B$ onto $A$, and
    $B$ generates the same vectors as $A$, so
    Lemma~\ref{thm:mpug_dressed_proportionality} gives a nonzero scalar
    relating them against long words.
\end{proof}

\begin{lemma}[Compatibility of the L-symbol with the anomaly]
    \label{thm:mpug_single_block_l_compatibility}
    \lean{MPOTensor.GroupFamily.ActionData.isCompatible_lSymbol}
    \leanok
    \uses{def:mpug_single_block_l_symbol, thm:mpug_single_block_l_symbol_spec,
        def:mpug_anomaly_three_cochain, def:mpug_l_symbols}
    In the setting of Lemma~\ref{thm:mpug_single_block_l_symbol_spec},
    (\ref{eq:mpug_l_compatibility}) for a one-point set of blocks holds for
    all $g,h,k\in G$:
    \begin{align}
        L_{g,hk}L_{h,k}
         & =\omega(g,h,k)\,L_{g,h}L_{gh,k}.
    \end{align}
\end{lemma}

\begin{proof}\leanok
    \uses{thm:mpug_single_block_l_symbol_spec,
        thm:mpug_action_tensor_reductions, thm:mpug_reduction_composition,
        thm:mpug_anomaly_three_cochain_spec, thm:mpug_dressed_proportionality}
    For the triple action $B=((O_gO_h)O_k)\cdot A$ and the two fusion trees
    $V^L$, $V^R$ of $(g,h,k)$ in
    Definition~\ref{def:mpug_anomaly_three_cochain}, the naturality and
    pentagon identities of $\mathsf a$ give five identities between left
    boundaries of reductions of $B$ onto $A$:
    \begin{align}
        \alpha_{g,h}(\mathbf 1\otimes V_k)\,\mathsf a^{-1}
         & =V_g(\mathbf 1\otimes\alpha_{h,k})\,\mathsf a^{-1}
        (\mathsf a^{-1}\otimes\mathbf 1),
        \label{eq:mpug_l_compat_first_boundary}               \\
        V_g(\mathbf 1\otimes\varphi_{h,k})\,\mathsf a^{-1}
        (\mathsf a^{-1}\otimes\mathbf 1)
         & =\alpha_{g,hk}
        \bigl((\mathbf 1\otimes F^<_{h,k})\,\mathsf a^{-1}\otimes\mathbf 1\bigr),
        \notag                                                \\
        \varphi_{g,h}(\mathbf 1\otimes V_k)\,\mathsf a^{-1}
         & =\alpha_{gh,k}
        \bigl((F^<_{g,h}\otimes\mathbf 1)\otimes\mathbf 1\bigr),
        \notag                                                \\
        \varphi_{gh,k}\bigl((F^<_{g,h}\otimes\mathbf 1)\otimes\mathbf 1\bigr)
         & =V_{ghk}(V^L\otimes\mathbf 1),
        \notag                                                \\
        \varphi_{g,hk}
        \bigl((\mathbf 1\otimes F^<_{h,k})\,\mathsf a^{-1}\otimes\mathbf 1\bigr)
         & =V_{ghk}(V^R\otimes\mathbf 1).
        \notag
    \end{align}
    Starting from $\alpha_{g,h}(\mathbf 1\otimes V_k)\mathsf a^{-1}B^{\mathbf w}$,
    replacing $\alpha_{g,h}$ by $L_{g,h}\varphi_{g,h}$, $\alpha_{gh,k}$ by
    $L_{gh,k}\varphi_{gh,k}$, and $V^L$ by $\omega(g,h,k)V^R$ along one
    path, or applying (\ref{eq:mpug_l_compat_first_boundary}) and replacing
    $\alpha_{h,k}$ by $L_{h,k}\varphi_{h,k}$ inside the action of $g$ and
    $\alpha_{g,hk}$ by $L_{g,hk}\varphi_{g,hk}$ along the other, uses only
    replacements that hold against long words by
    Lemma~\ref{thm:mpug_single_block_l_symbol_spec}, lifted through the
    action tensors by Lemma~\ref{thm:mpug_action_tensor_reductions} and
    pulled back along the reductions by
    Lemma~\ref{thm:mpug_dressed_proportionality}. For all long words,
    \begin{align}
        \alpha_{g,h}(\mathbf 1\otimes V_k)\,\mathsf a^{-1}B^{\mathbf w}
         & =L_{g,h}L_{gh,k}\,\omega(g,h,k)\,V_{ghk}(V^R\otimes\mathbf 1)
        B^{\mathbf w}
        =L_{h,k}L_{g,hk}\,V_{ghk}(V^R\otimes\mathbf 1)B^{\mathbf w}.
        \notag
    \end{align}
    The boundary $V_{ghk}(V^R\otimes\mathbf 1)$ belongs to a reduction onto
    the normal tensor $A$, so it is nonzero against arbitrarily long words,
    and the two scalars agree.
\end{proof}

\begin{theorem}[An anomaly obstructs invariant normal states]
    \label{thm:mpug_anomaly_obstruction}
    \lean{MPOTensor.GroupFamily.IsNormalRepresentation.isTrivialGaugeClass_omega_of_invariant,
        MPOTensor.GroupFamily.IsNormalRepresentation.not_exists_invariant_of_not_isTrivialGaugeClass}
    \leanok
    \uses{def:mpug_anomaly_three_cochain, def:mpug_trivial_gauge_class,
        def:mpv_state}
    Let $\{O_g\}$ be matrix product operator tensors with normal
    doubled-index tensors and $O_gO_h=O_{gh}$ on every nonempty chain. If a
    normal matrix product state tensor of positive bond dimension satisfies
    $O_g\ket{V^{(N)}(A)}=\ket{V^{(N)}(A)}$ for every $g$ and every $N\geq1$, then
    the anomaly three-cocycle has trivial gauge class. Equivalently, as in
    \cite[line~1545]{FrancoRubio2025MPUGauging} and
    \cite[line~738]{GarreRubio2023MPOSym} with normal in place of injective
    and without unitarity or simplicity of the operators, a group with
    nontrivial anomaly has no invariant normal matrix product state.
\end{theorem}

\begin{proof}\leanok
    \uses{thm:mpug_single_block_l_symbol_spec,
        thm:mpug_single_block_l_compatibility,
        thm:mpug_fixed_block_trivializes}
    Action tensors exist by
    Lemma~\ref{thm:mpug_single_block_l_symbol_spec}, so
    Lemma~\ref{thm:mpug_single_block_l_compatibility} applies, and Lemma~\ref{thm:mpug_fixed_block_trivializes} for the
    single block gives $\omega=dL$.
\end{proof}

\subsection{Blocks permuted by the group}

Let $\{O_g\}$ be as in Definition~\ref{def:mpug_anomaly_three_cochain}, and
let $G$ act on a set $\mathsf X$. Following
\cite[line~1064]{GarreRubio2023MPOSym}, consider normal matrix product
state tensors $A_x$, $x\in\mathsf X$, of positive bond dimensions, possibly
different for different $x$, with
$O_g\ket{V^{(N)}(A_x)}=\ket{V^{(N)}(A_{gx})}$ for every $g$, $x$ and
$N\geq1$.

\begin{definition}[L-symbols of permuted blocks]
    \label{def:mpug_permuted_block_l_symbol}
    \lean{MPOTensor.GroupFamily.CarriesMPV,
        MPOTensor.GroupFamily.castIndex,
        MPOTensor.GroupFamily.BlockActionData,
        MPOTensor.GroupFamily.BlockActionData.actV,
        MPOTensor.GroupFamily.BlockActionData.fuseV,
        MPOTensor.GroupFamily.BlockActionData.lSymbol}
    \leanok
    \uses{def:mpug_anomaly_three_cochain, def:mpug_l_symbols,
        def:mpug_single_block_l_symbol}
    Action tensors for the family $(A_x)$ are, for every $g$ and $x$, a
    reduction $(V_{g,x},W_{g,x})$ of $O_g\cdot A_x$ onto $A_{gx}$. The stacked
    action $B=(O_gO_h)\cdot A_x$ then reduces onto $A_{g(hx)}=A_{(gh)x}$ along
    two trees, with left boundaries
    \begin{align}
        V_{g,hx}(\mathbf 1\otimes V_{h,x})\,\mathsf a^{-1}, &  &
        V_{gh,x}(F^<_{g,h}\otimes\mathbf 1),
        \notag
    \end{align}
    and $L^x_{g,h}$ is a nonzero scalar relating them against all long
    words of $B$, or $1$ if there is none, as in
    Definition~\ref{def:mpug_single_block_l_symbol}.
\end{definition}

\begin{lemma}[Characterization of the L-symbols of permuted blocks]
    \label{thm:mpug_permuted_block_l_symbol_spec}
    \lean{MPOTensor.GroupFamily.nonempty_blockActionData,
        MPOTensor.GroupFamily.BlockActionData.isDressedProportional_lSymbol,
        MPOTensor.GroupFamily.BlockActionData.eq_lSymbol_of_isDressedProportional,
        MPOTensor.GroupFamily.actTensor_idTensor_apply,
        MPOTensor.GroupFamily.isReduction_actTensor_idTensor}
    \leanok
    \uses{def:mpug_permuted_block_l_symbol}
    Action tensors for the family $(A_x)$ of
    Definition~\ref{def:mpug_permuted_block_l_symbol} exist; for the bond-one
    identity tensor they are the identification of the bond spaces. For
    every choice of them, $L^x_{g,h}$ relates the two left boundaries against
    all long words of $(O_gO_h)\cdot A_x$, and it is the only scalar that
    does.
\end{lemma}

\begin{proof}\leanok
    \uses{cor:asym_action_tensors, thm:mpug_action_tensor_reductions,
        thm:mpug_reduction_composition, thm:mpug_dressed_proportionality}
    As in Lemma~\ref{thm:mpug_single_block_l_symbol_spec}: existence is
    Corollary~\ref{cor:asym_action_tensors} applied to $O_g\cdot A_x$ and
    $A_{gx}$, which generate the same vectors. By
    Lemmas~\ref{thm:mpug_action_tensor_reductions}
    and~\ref{thm:mpug_reduction_composition} both boundaries are left
    boundaries of reductions of $(O_gO_h)\cdot A_x$ onto the normal tensor
    $A_{(gh)x}$, so Lemma~\ref{thm:mpug_dressed_proportionality} gives a
    nonzero scalar relating them against long words. The boundary of two
    successive actions is nonzero against arbitrarily long words, since it is
    the left boundary of a reduction onto a normal tensor, so the scalar is
    unique.
\end{proof}

\begin{lemma}[Blocking the action tensors]
    \label{lem:mpug_action_blocking}
    \lean{MPOTensor.blockTensor_actTensor,
        MPOTensor.GroupFamily.CarriesMPV.blockTensor,
        MPOTensor.GroupFamily.BlockActionData.block}
    \leanok
    \uses{def:mpug_permuted_block_l_symbol, def:mpug_action_tensor}
    For an operator tensor $T$ and a state tensor $A$, physical blocking
    satisfies $(T\cdot A)^{[L]}=T^{[L]}\cdot A^{[L]}$. A reduction
    $(V,W)$ from $T\cdot A$ onto $C$ remains a reduction from
    $T^{[L]}\cdot A^{[L]}$ onto $C^{[L]}$, with the same matrices.
    For $L>0$, the periodic action law
    $O_N(T)\ket{V^{(N)}(A)}=\ket{V^{(N)}(C)}$ at all positive lengths
    is also preserved by blocking.
\end{lemma}

\begin{proof}\leanok
    \uses{def:mpug_action_tensor, thm:mps_reduction_exterior_buffer}
    Decode a blocked input index as a word of length $L$. Reindexing
    the contracted input sum gives
    \begin{align}
        (T\cdot A)^{[L],I}
         & =\sum_J T^{[L],IJ}\otimes A^{[L],J}.
        \label{eq:mpug_action_blocking}
    \end{align}
    A word of blocked letters is the corresponding flattened word of
    original letters. The reduction equations and the periodic action law
    therefore apply at the flattened lengths.
\end{proof}

\begin{definition}[A common action buffer]
    \label{def:mpug_action_buffer}
    \lean{MPOTensor.GroupFamily.BlockActionData.nilpotencyLength,
        MPOTensor.GroupFamily.BlockActionData.HasExteriorBufferLength,
        MPOTensor.GroupFamily.BlockActionData.commonBufferLength,
        MPOTensor.GroupFamily.BlockActionData.commonBufferLength_pos,
        MPOTensor.GroupFamily.BlockActionData.nilpotencyLength_le_commonBufferLength_add_one,
        MPOTensor.GroupFamily.BlockActionData.commonBufferLength_of_isEmpty}
    \leanok
    \uses{def:mpug_permuted_block_l_symbol, thm:mps_reduction_exterior_buffer}
    Fix action reductions $(V_{g,x},W_{g,x})$ for finite $G$ and
    $\mathsf X$, and put $B_{g,x}=O_g\cdot A_x$. Write $N_0(g,x)$ for
    their residual nilpotency lengths. Their common buffer length is
    $L=\max\{1,\max_{g,x}(N_0(g,x)-1)\}$, where the inner maximum is
    zero for the empty family. Thus $L>0$ and $N_0(g,x)\leq L+1$;
    when $\mathsf X$ is empty, $L=1$.
\end{definition}

\begin{lemma}[Common blocking of action reductions]
    \label{lem:mpug_action_common_blocking}
    \lean{MPOTensor.GroupFamily.BlockActionData.interior,
        MPOTensor.GroupFamily.BlockActionData.exterior,
        MPOTensor.GroupFamily.BlockActionData.block_hasExteriorBufferLength_one,
        MPOTensor.GroupFamily.BlockActionData.block_commonBufferLength_hasExteriorBufferLength_one,
        MPOTensor.GroupFamily.exists_block_blockActionData_one,
        MPOTensor.GroupFamily.exists_block_blockActionData_one_of_isPretransitive}
    \leanok
    \uses{def:mpug_action_buffer, lem:mpug_action_blocking,
        thm:mpug_permuted_block_l_symbol_spec}
    Suppose $G$ and $\mathsf X$ are finite, the tensors $A_x$ are normal
    of positive, possibly different, bond dimensions, and the periodic
    action law holds at every positive length. There is one positive
    length $L$ for which all the blocked action tensors
    $B_{g,x}=O_g^{[L]}\cdot A_x^{[L]}$ have reductions
    $(V_{g,x},W_{g,x})$ onto $C_{g,x}=A_{gx}^{[L]}$ satisfying
    \begin{align}
        V_{g,x}B_{g,x}^{\mathbf w}W_{g,x}
         & =C_{g,x}^{\mathbf w},
        \label{eq:mpug_action_interior}            \\
        B_{g,x}^{\mathbf p}W_{g,x}C_{g,x}^{\mathbf c}
        V_{g,x}B_{g,x}^{\mathbf q}
         & =B_{g,x}^{\mathbf p\mathbf c\mathbf q}.
        \label{eq:mpug_action_exterior}
    \end{align}
    The interior word $\mathbf w$ may be empty. In the exterior identity
    $\mathbf c$ is nonempty and $|\mathbf p|,|\mathbf q|\geq1$.
    These are the interior and exterior equations of
    \cite[lines~1782--1873]{FrancoRubio2025MPUGauging}, with
    $A^\Gamma_{g,x}=V_{g,x}$ and
    $A^{\text{\reflectbox{$\Gamma$}}}_{g,x}=W_{g,x}$.
    The periodic action hypothesis is explicit, as in
    \cite[line~1064]{GarreRubio2023MPOSym}; see also
    \cite{gap:fbc25_mpu_action_stabilizer_source_audit}.

    In particular, for a transitive action of finite $G$ on $\mathsf X$
    and injective $A_x$, the label set is automatically finite. If
    $\{O_g\}$ is a simple injective MPU representation, this same blocking
    preserves that property, the injectivity of every $A_x$, and the
    periodic action law.
\end{lemma}

\begin{proof}\leanok
    \uses{def:mpug_action_buffer, lem:mpug_action_blocking,
        thm:mpug_permuted_block_l_symbol_spec, thm:mps_reduction_exterior_buffer}
    Choose the action reductions and the common length of
    Definition~\ref{def:mpug_action_buffer}. Their residual bounds follow
    from equality of the positive-length periodic vectors. Since
    $N_0(g,x)-1\leq L$, one exterior block supplies the required buffer
    for every pair. The interior identity holds on all words by the
    reduction equations. If $\mathsf X$ is nonempty and the action is
    transitive, the map $g\mapsto gx_0$ is onto, so $\mathsf X$ is finite;
    the empty case needs no choice of $x_0$. Positive blocking preserves
    injectivity and the MPU representation laws.
\end{proof}

\begin{lemma}[Reciprocal scalar action gauge]
    \label{lem:mpug_action_scalar_gauge}
    \lean{MPOTensor.GroupFamily.BlockActionData.smul,
        MPOTensor.GroupFamily.BlockActionData.hasExteriorBufferLength_smul,
        MPOTensor.GroupFamily.BlockActionData.block_smul}
    \leanok
    \uses{lem:mpug_action_common_blocking}
    For any nonzero scalars $\gamma_{g,x}$ the replacements
    $W_{g,x}\mapsto\gamma_{g,x}W_{g,x}$ and
    $V_{g,x}\mapsto\gamma_{g,x}^{-1}V_{g,x}$ preserve
    (\ref{eq:mpug_action_interior}) and
    (\ref{eq:mpug_action_exterior}) with the same exterior buffer length.
    They commute with physical blocking. This is the scalar action in
    \cite[lines~1871--1873]{FrancoRubio2025MPUGauging}.
\end{lemma}

\begin{proof}\leanok
    The two scalars cancel inside $V_{g,x}B_{g,x}^{\mathbf w}W_{g,x}$
    and $W_{g,x}C_{g,x}^{\mathbf c}V_{g,x}$. Blocking changes the physical
    tensors and leaves the reduction matrices unchanged.
\end{proof}


\begin{lemma}[Compatibility for permuted blocks]
    \label{thm:mpug_permuted_block_l_compatibility}
    \lean{MPOTensor.GroupFamily.BlockActionData.isCompatible_lSymbol,
        MPOTensor.GroupFamily.IsNormalRepresentation.isTrivialGaugeClass_comap_omega_of_fixed}
    \leanok
    \uses{def:mpug_permuted_block_l_symbol, thm:mpug_anomaly_three_cochain_spec,
        thm:mpug_permuted_block_l_symbol_spec,
        thm:mpug_dressed_proportionality, thm:mpug_fixed_block_trivializes}
    For all $x\in\mathsf X$ and $g,h,k\in G$,
    \begin{align}
        L^x_{g,hk}L^x_{h,k}
         & =\omega(g,h,k)\,L^{kx}_{g,h}L^x_{gh,k},
    \end{align}
    which is (\ref{eq:mpug_l_compatibility}). Consequently, as in
    \cite[lines~740--745]{GarreRubio2023MPOSym}, the restriction of
    $\omega$ to any subgroup fixing one block has trivial gauge class.
\end{lemma}

\begin{proof}\leanok
    \uses{thm:mpug_permuted_block_l_symbol_spec,
        thm:mpug_anomaly_three_cochain_spec,
        thm:mpug_fixed_block_trivializes}
    Write
    \begin{align}
        \alpha^x_{g,h}  & =V_{g,hx}(\mathbf 1\otimes V_{h,x})\,\mathsf a^{-1}, &
        \varphi^x_{g,h} & =V_{gh,x}(F^<_{g,h}\otimes\mathbf 1)
        \notag
    \end{align}
    for the two left boundaries of Definition~\ref{def:mpug_permuted_block_l_symbol},
    so that $\alpha^x_{g,h}=L^x_{g,h}\varphi^x_{g,h}$ against long words of
    $(O_gO_h)\cdot A_x$ by Lemma~\ref{thm:mpug_permuted_block_l_symbol_spec}.
    The block labels $h(kx)$ and $(hk)x$ are equal, and so are $g(h(kx))$ and
    $(ghk)x$, and the bond spaces of blocks with equal labels are identified,
    so that $V_{g,h(kx)}=V_{g,(hk)x}$. With this identification, the
    naturality and
    pentagon identities of $\mathsf a$ give, for $B=((O_gO_h)O_k)\cdot A_x$
    and the fusion trees $V^L$, $V^R$ of $(g,h,k)$, the five identities
    \begin{align}
        \alpha^{kx}_{g,h}(\mathbf 1\otimes V_{k,x})\,\mathsf a^{-1}
         & =V_{g,(hk)x}(\mathbf 1\otimes\alpha^x_{h,k})\,\mathsf a^{-1}
        (\mathsf a^{-1}\otimes\mathbf 1),
        \label{eq:mpug_permuted_l_compat_first_boundary}                \\
        V_{g,(hk)x}(\mathbf 1\otimes\varphi^x_{h,k})\,\mathsf a^{-1}
        (\mathsf a^{-1}\otimes\mathbf 1)
         & =\alpha^x_{g,hk}
        \bigl((\mathbf 1\otimes F^<_{h,k})\,\mathsf a^{-1}\otimes\mathbf 1\bigr),
        \notag                                                          \\
        \varphi^{kx}_{g,h}(\mathbf 1\otimes V_{k,x})\,\mathsf a^{-1}
         & =\alpha^x_{gh,k}
        \bigl((F^<_{g,h}\otimes\mathbf 1)\otimes\mathbf 1\bigr),
        \notag                                                          \\
        \varphi^x_{gh,k}\bigl((F^<_{g,h}\otimes\mathbf 1)\otimes\mathbf 1\bigr)
         & =V_{ghk,x}(V^L\otimes\mathbf 1),
        \notag                                                          \\
        \varphi^x_{g,hk}
        \bigl((\mathbf 1\otimes F^<_{h,k})\,\mathsf a^{-1}\otimes\mathbf 1\bigr)
         & =V_{ghk,x}(V^R\otimes\mathbf 1).
        \notag
    \end{align}
    The first factor $\alpha^{kx}_{g,h}$ acts on the block $kx$ reached by
    $O_k$, which is why $L^{kx}_{g,h}$ appears. Replacing, as in the proof of
    Lemma~\ref{thm:mpug_single_block_l_compatibility}, $\alpha^{kx}_{g,h}$,
    $\alpha^x_{gh,k}$ and $V^L$ along one route, and
    (\ref{eq:mpug_permuted_l_compat_first_boundary}), $\alpha^x_{h,k}$ and
    $\alpha^x_{g,hk}$ along the other, gives for all long words
    \begin{align}
        \alpha^{kx}_{g,h}(\mathbf 1\otimes V_{k,x})\,\mathsf a^{-1}B^{\mathbf w}
         & =L^{kx}_{g,h}L^x_{gh,k}\,\omega(g,h,k)\,
        V_{ghk,x}(V^R\otimes\mathbf 1)B^{\mathbf w}
        \notag                                                                \\
         & =L^x_{h,k}L^x_{g,hk}\,V_{ghk,x}(V^R\otimes\mathbf 1)B^{\mathbf w}.
        \notag
    \end{align}
    The boundary $V_{ghk,x}(V^R\otimes\mathbf 1)$ belongs to a reduction onto
    the normal tensor $A_{(ghk)x}$, so it is nonzero against arbitrarily long
    words, and the two scalars agree. The second statement is
    Lemma~\ref{thm:mpug_fixed_block_trivializes}.
\end{proof}

\begin{lemma}[The anomaly at an involution]
    \label{thm:mpug_involution_anomaly_ratio}
    \lean{TNLean.Algebra.LSymbol.eq_div_of_isCompatible_of_mul_self_eq_one}
    \leanok
    \uses{def:mpug_l_symbols}
    If $L$ is compatible with $\omega$ and $g^2=e$, then for every
    $x\in\mathsf X$ the formula of
    \cite[line~5121]{FrancoRubio2025MPUGauging} holds:
    \begin{align}
        \omega(g,g,g)
         & =\frac{L^x_{g,g}L^x_{g,e}}{L^{gx}_{g,g}L^x_{e,g}}.
    \end{align}
\end{lemma}

\begin{proof}\leanok
    Setting $h=k=g$ in (\ref{eq:mpug_l_compatibility}) and using $gg=e$
    gives the ratio.
\end{proof}

\begin{lemma}[The involution ratio from the self-pair and identity L-symbols]
    \label{thm:mpug_involution_anomaly_identity_ratio}
    \lean{TNLean.Algebra.LSymbol.div_eq_mul_of_isCompatible_of_mul_self_eq_one}
    \leanok
    \uses{def:mpug_l_symbols}
    If $L$ is compatible with $\omega$ and $g^2=e$, then for every
    $x\in\mathsf X$
    \begin{align}
        \frac{L^x_{g,g}L^x_{e,e}}{L^{gx}_{g,g}L^{gx}_{e,e}}
         & =\omega(g,g,g)\,\omega(g,e,g).
        \notag
    \end{align}
    For L-symbols and cocycles that are trivial when an argument is the
    identity, this recovers the relation $L_A/L_B=\omega$ of
    \cite[Section~III.C]{GarreRubioSchuch2024DomainWall} in the form
    $L^x_{g,g}/L^{gx}_{g,g}=\omega(g,g,g)$.
\end{lemma}

\begin{proof}\leanok
    \uses{thm:mpug_involution_anomaly_ratio}
    Compatibility at $(g,e,g)$ gives $L^x_{e,g}=\omega(g,e,g)L^{gx}_{g,e}$,
    and at $(g,e,e)$ it gives $L^z_{e,e}=\omega(g,e,e)L^z_{g,e}$ for $z=x$ and
    $z=gx$. Substituting into the formula of
    Lemma~\ref{thm:mpug_involution_anomaly_ratio} gives the claim.
\end{proof}

\begin{lemma}[Gauge law for the L-symbols of permuted blocks]
    \label{lem:mpug_permuted_block_l_gauge}
    \lean{MPOTensor.GroupFamily.FusionData.rescale,
        MPOTensor.GroupFamily.FusionData.isDressedProportional_rescale,
        MPOTensor.GroupFamily.FusionData.relativeScalar_rescale,
        MPOTensor.GroupFamily.FusionData.omega_rescale,
        MPOTensor.GroupFamily.BlockActionData.relativeScalar,
        MPOTensor.GroupFamily.BlockActionData.isDressedProportional_relativeScalar,
        MPOTensor.GroupFamily.BlockActionData.rescale,
        MPOTensor.GroupFamily.BlockActionData.relativeScalar_rescale,
        MPOTensor.GroupFamily.BlockActionData.lSymbol_eq_gauge,
        MPOTensor.GroupFamily.BlockActionData.lSymbol_rescale}
    \leanok
    \uses{def:mpug_permuted_block_l_symbol,
        thm:mpug_permuted_block_l_symbol_spec,
        thm:mpug_anomaly_three_cocycle_gauge, def:mpug_l_gauge,
        thm:mpug_dressed_proportionality}
    Keep the setting of Definition~\ref{def:mpug_permuted_block_l_symbol}
    and let the doubled-index tensors be normal. Take two choices of fusion
    tensors and two choices of action tensors, with $L$ and $L'$ the
    corresponding $L$-symbols. Let $\beta_{g,h}$ be the nonzero scalar with
    $F^<_{g,h}B^{\mathbf w}=\beta_{g,h}F'^<_{g,h}B^{\mathbf w}$ for all long
    words of $B=O_gO_h$, as in
    Lemma~\ref{thm:mpug_anomaly_three_cocycle_gauge}, and let $\gamma^x_g$ be
    the nonzero scalar with
    $V_{g,x}C^{\mathbf w}=\gamma^x_gV'_{g,x}C^{\mathbf w}$ for all long words
    of $C=O_g\cdot A_x$. Then
    \begin{align}
        L'^x_{g,h}
         & =\frac{\gamma^x_{gh}\beta_{g,h}}{\gamma^{hx}_g\gamma^x_h}L^x_{g,h},
        \notag
    \end{align}
    which is the joint gauge (\ref{eq:mpug_l_gauge}). Conversely,
    multiplying $F^>_{g,h}$ by $\beta_{g,h}$ and $F^<_{g,h}$ by
    $\beta_{g,h}^{-1}$, and $W_{g,x}$ by $\gamma^x_g$ and $V_{g,x}$ by
    $(\gamma^x_g)^{-1}$, gives fusion and action tensors whose comparison
    scalars with the original ones are exactly $\beta$ and $\gamma$, so every
    joint scalar gauge of the $L$-symbols is realized by tensors, and the
    rescaled fusion tensors have anomaly three-cocycle
    $\beta\mathbin{\triangleright}\omega$.

    This is the gauge law \cite[line~1918]{FrancoRubio2025MPUGauging}, with
    the scalars of
    \cite[lines~1500--1504, 1871--1873]{FrancoRubio2025MPUGauging}.
    As recorded in
    \cite{gap:fbc25_boundary_dressed_l_symbols}, all three families of scalars
    relate the tensors against long words rather than as bare matrices,
    because the uniqueness of reductions onto a normal tensor holds only in
    that form.
\end{lemma}

\begin{proof}\leanok
    \uses{thm:mpug_permuted_block_l_symbol_spec,
        thm:mpug_dressed_proportionality}
    As in the proof of Lemma~\ref{thm:mpug_permuted_block_l_compatibility},
    write $\alpha^x_{g,h}$ and $\varphi^x_{g,h}$ for the two left boundaries
    of Definition~\ref{def:mpug_permuted_block_l_symbol}, and
    $\alpha'^x_{g,h}$, $\varphi'^x_{g,h}$ for those of the primed choice.
    Comparing the action tensors of $h$ on the block $x$ and lifting through
    the operator $O_g$, then pulling the comparison for $g$ on the block $hx$
    back along the reduction by the primed action of $h$, and comparing the
    fusion tensors in the same way, gives for all long words of
    $C=(O_gO_h)\cdot A_x$
    \begin{align}
        \alpha^x_{g,h}C^{\mathbf w}
         & =\gamma^{hx}_g\gamma^x_h\,\alpha'^x_{g,h}C^{\mathbf w},   &
        \varphi^x_{g,h}C^{\mathbf w}
         & =\beta_{g,h}\gamma^x_{gh}\,\varphi'^x_{g,h}C^{\mathbf w}.
        \label{eq:mpug_permuted_l_gauge_boundaries}
    \end{align}
    Since $\alpha=L\varphi$ and $\alpha'=L'\varphi'$ against long words, the
    first boundary of the original choice equals both
    $\gamma^{hx}_g\gamma^x_hL'^x_{g,h}\varphi'^x_{g,h}$ and
    $L^x_{g,h}\beta_{g,h}\gamma^x_{gh}\varphi'^x_{g,h}$ against long words,
    and $\varphi'^x_{g,h}$ is nonzero against arbitrarily long words because
    it belongs to a reduction onto a normal tensor. For rescaled tensors the left boundaries
    differ by the rescaling factor exactly, so the comparison scalars are the
    rescaling factors, and the statement about the anomaly is
    Lemma~\ref{thm:mpug_anomaly_three_cocycle_gauge}.
\end{proof}

\begin{lemma}[A single block is trivial in some gauge]
    \label{lem:mpug_single_block_gauge_trivial}
    \lean{MPOTensor.GroupFamily.BlockActionData.exists_rescale_omega_eq_one_lSymbol_eq_one}
    \leanok
    \uses{lem:mpug_permuted_block_l_gauge,
        thm:mpug_permuted_block_l_compatibility}
    In the setting of Lemma~\ref{lem:mpug_permuted_block_l_gauge}, if the set
    $\mathsf X$ has exactly one element, some rescaling of the fusion and
    action tensors gives $\omega(g,h,k)=1$ and $L_{g,h}=1$ for all
    $g,h,k\in G$, as stated in
    \cite[lines~1925--1926]{FrancoRubio2025MPUGauging}.
\end{lemma}

\begin{proof}\leanok
    \uses{lem:mpug_permuted_block_l_gauge,
        thm:mpug_permuted_block_l_compatibility}
    Rescale the fusion tensors by $\beta_{g,h}=L_{g,h}^{-1}$ and leave the
    action tensors unchanged. By Lemma~\ref{lem:mpug_permuted_block_l_gauge}
    the new $L$-symbols are $\beta_{g,h}L_{g,h}=1$, and compatibility,
    Lemma~\ref{thm:mpug_permuted_block_l_compatibility}, for the new tensors
    then reads $1=\omega'(g,h,k)$.
\end{proof}

\subsection{Domain walls between permuted blocks}

Keep the setting of the previous subsection. A domain wall between the
blocks $x$ and $y$ is a family $e=(e^i)_{i}$ of $D_x\times D_y$ matrices,
placed between a region of $A_x$ and a region of $A_y$
\cite[Section~III.B]{GarreRubioSchuch2024DomainWall}. For an operator tensor
$T$ write $(T\cdot e)^i=\sum_j T^{ij}\otimes e^j$.

\begin{definition}[Local action on domain walls]
    \label{def:mpug_domain_wall_action}
    \lean{MPOTensor.actRect,
        MPOTensor.physAct,
        MPOTensor.wallChain,
        MPOTensor.GroupFamily.BlockActionData.IsDomainWallAction}
    \leanok
    \uses{def:mpug_permuted_block_l_symbol}
    Let $e$ be a domain wall between $x$ and $y$, and $e'$ a domain wall
    between $gx$ and $gy$. The element $g$ carries $e$ to $e'$ with phase $c$
    if $e\neq0$, $e'\neq0$, $c\neq0$, and there is $N$ such that
    \begin{align}
        V_{g,x}\,(O_g\cdot A_x)^{\mathbf u}\,(O_g\cdot e)^i\,
        (O_g\cdot A_y)^{\mathbf v}\,W_{g,y}
         & =c\,A_{gx}^{\mathbf u}\,{e'}^{i}\,A_{gy}^{\mathbf v}
        \notag
    \end{align}
    for every letter $i$ and all words $\mathbf u,\mathbf v$ of length at
    least $N$. This is the local relation of
    \cite[Section~III.B]{GarreRubioSchuch2024DomainWall}, drawn there at the
    renormalization fixed point with one site on each side of the wall; the
    words $\mathbf u,\mathbf v$ are the blocked sites of the discussion away
    from the fixed point \cite{gap:gs24_domain_wall_nondegenerate}.
\end{definition}

\begin{lemma}[Fractionalization of the symmetry on domain walls]
    \label{thm:mpug_domain_wall_pentagon}
    \lean{MPOTensor.GroupFamily.BlockActionData.IsDomainWallAction.mul,
        MPOTensor.GroupFamily.BlockActionData.IsDomainWallAction.mul_eq,
        MPOTensor.GroupFamily.BlockActionData.IsDomainWallAction.phase_eq,
        MPOTensor.GroupFamily.BlockActionData.actW,
        MPOTensor.GroupFamily.BlockActionData.fuseW,
        MPOTensor.GroupFamily.BlockActionData.exists_evalWord_mul_actW_eq,
        MPOTensor.physAct_mulTensor,
        MPOTensor.physAct_kronId,
        MPOTensor.physAct_wallChain,
        MPOTensor.exists_evalWord_mul_mul_evalWord_ne_zero}
    \leanok
    \uses{def:mpug_domain_wall_action, def:mpug_permuted_block_l_symbol}
    Suppose $h$ carries the domain wall $e_1$ between $x$ and $y$ to $e_2$
    with phase $B^h_{x,y}$ and $g$ carries $e_2$ to $e_3$ with phase
    $B^g_{hx,hy}$. Then $gh$ carries $e_1$ to $e_3$ with phase
    $(L^y_{g,h}/L^x_{g,h})\,B^g_{hx,hy}B^h_{x,y}$. Consequently, as in
    \cite[Section~IV.B]{GarreRubioSchuch2024DomainWall}, if $gh$
    carries $e_1$ to $e_3$ with phase $B^{gh}_{x,y}$, then
    \begin{align}
        B^g_{hx,hy}\,B^h_{x,y}
         & =\frac{L^x_{g,h}}{L^y_{g,h}}\,B^{gh}_{x,y}.
        \notag
    \end{align}
\end{lemma}

\begin{proof}\leanok
    \uses{thm:mpug_permuted_block_l_symbol_spec, thm:mpug_dressed_proportionality}
    Acting with $O_g$ on the local relation of $h$, which holds for every
    configuration of a chain with long regions, passes the boundary matrices
    $V_{h,x}$, $W_{h,y}$ through as $\mathbf 1\otimes V_{h,x}$ and
    $\mathbf 1\otimes W_{h,y}$, while the action of $O_g$ on the action of
    $O_h$ is the action of $O_gO_h$ up to the bond associator. The local
    relation of $g$ then gives
    \begin{align}
        V_{g,hx}(\mathbf 1\otimes V_{h,x})\,\mathsf a^{-1}\,
        (O_gO_h\cdot\mathsf C)\,\mathsf a\,(\mathbf 1\otimes W_{h,y})\,W_{g,hy}
         & =B^g_{hx,hy}B^h_{x,y}\,\mathsf C_3
        \notag
    \end{align}
    for the chain $\mathsf C=A_x^{\mathbf u}e_1^iA_y^{\mathbf v}$ and the chain
    $\mathsf C_3$ with $e_3$. Against long words the left boundary equals
    $L^x_{g,h}$ times the left boundary through the fusion tensor
    $F^<_{g,h}$; by the two-sided comparison of reductions onto a normal
    tensor, the right boundary equals $(L^y_{g,h})^{-1}$ times the right
    boundary through $F^>_{g,h}$. The fusion tensors reduce
    $O_gO_h\cdot\mathsf C$ to $O_{gh}\cdot\mathsf C$, which gives the local
    action of $gh$ with the stated phase. Since $A_{ghx}$ and $A_{ghy}$ are
    normal and $e_3\neq0$, some chain $\mathsf C_3$ with long regions is
    nonzero, so the phase of a local action is unique, and the two phases of
    $gh$ agree.
\end{proof}

\begin{lemma}[The symmetry on states with two domain walls]
    \label{thm:mpug_domain_wall_state}
    \lean{MPOTensor.twoWallChain,
        MPOTensor.twoWallMPV,
        MPOTensor.mpo_mulVec_trace,
        MPOTensor.GroupFamily.BlockActionData.IsDomainWallAction.mpo_mulVec_twoWallMPV}
    \leanok
    \uses{def:mpug_domain_wall_action}
    Let $\ket{\psi(x\text{-}y\text{-}x)}$ be the periodic state with an $A_x$
    region, a domain wall $e$ between $x$ and $y$, an $A_y$ region, a domain
    wall $f$ between $y$ and $x$, and a second $A_x$ region. If $g$ carries
    $e$ to $e'$ with phase $c$ and $f$ to $f'$ with phase $c'$, then
    \begin{align}
        O_g\ket{\psi(x\text{-}y\text{-}x)}
         & =c\,c'\,\ket{\psi(gx\text{-}gy\text{-}gx)}
        \notag
    \end{align}
    whenever the first $A_x$ region and the $A_y$ region are long enough.
    For $\Z_2$, as in \cite[Section~III.B]{GarreRubioSchuch2024DomainWall},
    this reduces to $U\ket{\psi(A\text{-}B\text{-}A)}
        =c_{AB}c_{BA}\ket{\psi(B\text{-}A\text{-}B)}$.
\end{lemma}

\begin{proof}\leanok
    \uses{thm:mps_reduction_residual_dimension_bound,
        thm:mps_reduction_exterior_buffer}
    The coefficient of $O_g\ket{\psi}$ is the trace of the acted chain.
    Inserting the reduced block $W_{g,z}A_{gz}^{\mathbf c}V_{g,z}$ into each
    of the two acted regions leaves long words unchanged by the residual
    nilpotency of a reduction \cite{Molnar2018SemiInjective}. By cyclicity
    of the trace the chain splits into the two local relations, with phases
    $c$ and $c'$, and the remaining trace is the coefficient of
    $\ket{\psi(gx\text{-}gy\text{-}gx)}$.
\end{proof}

\begin{lemma}[A symmetry string over a pair of domain walls]
    \label{thm:mpug_domain_wall_pair}
    \lean{MPOTensor.GroupFamily.BlockActionData.IsDomainWallAction.pair}
    \leanok
    \uses{def:mpug_domain_wall_action}
    If $g$ carries the domain wall $e$ between $x$ and $y$ to $e'$ with phase
    $c$, and the domain wall $f$ between $y$ and $z$ to $f'$ with phase $c'$,
    then for all words $\mathbf u,\mathbf v,\mathbf w$ longer than a fixed
    buffer
    \begin{align}
        V_{g,x}\,(O_g\cdot A_x)^{\mathbf u}(O_g\cdot e)^i
        (O_g\cdot A_y)^{\mathbf v}(O_g\cdot f)^j(O_g\cdot A_z)^{\mathbf w}\,
        W_{g,z}
         & =c\,c'\,A_{gx}^{\mathbf u}{e'}^{i}A_{gy}^{\mathbf v}{f'}^{j}
        A_{gz}^{\mathbf w}.
        \notag
    \end{align}
    In the comparison of two overlapping strings in
    \cite[Section~III.F]{GarreRubioSchuch2024DomainWall}, this is the
    tensor-level phase acquired by a symmetry string passing over a pair of
    domain walls.
\end{lemma}

\begin{proof}\leanok
    \uses{thm:mps_reduction_residual_dimension_bound,
        thm:mps_reduction_exterior_buffer}
    Splitting the middle word as $\mathbf p\mathbf c\mathbf q$ with long
    outer parts and inserting the reduced block
    $W_{g,y}A_{gy}^{\mathbf c}V_{g,y}$ leaves the word unchanged by the
    residual nilpotency of a reduction \cite{Molnar2018SemiInjective}. The
    two resulting factors are the local relations of $e$ and $f$.
\end{proof}

\begin{definition}[Domain-wall string operators]
    \label{def:mpug_domain_wall_string}
    \lean{MPOTensor.physPairing,
        MPOTensor.leftEndpoint,
        MPOTensor.rightEndpoint,
        MPOTensor.leftAct,
        MPOTensor.rightAct,
        MPOTensor.wallSplit,
        MPOTensor.wallConfig,
        MPOTensor.wallConfig_comp_cast,
        MPOTensor.ofFn_wallConfig,
        MPOTensor.stringKernel,
        MPOTensor.stringOperator,
        MPOTensor.GroupFamily.IsSeparatingLeftInverse,
        MPOTensor.GroupFamily.BlockActionData.wallLeftEndpoint,
        MPOTensor.GroupFamily.BlockActionData.wallRightEndpoint,
        MPOTensor.GroupFamily.BlockActionData.wallString}
    \leanok
    \uses{def:mpug_domain_wall_action}
    Let $g$ exchange the blocks $A=A_x$ and $B=A_y$, with left and right action
    tensors $V_B=V_{g,x}$, $W_B=W_{g,x}$ on $A$ and $V_A=V_{g,y}$,
    $W_A=W_{g,y}$ on $B$, and let $e_{AB}$ and $e_{BA}$ be domain walls. Let
    $\hat A$ and $\hat B$ be left inverses of $A$ and $B$ with separated
    supports,
    \begin{align}
        \sum_\sigma \hat A^\sigma_{\alpha\beta}A^\sigma_{\gamma\delta}
         & =\delta_{\alpha\gamma}\delta_{\beta\delta},
         &
        \sum_\sigma \hat A^\sigma_{\alpha\beta}B^\sigma_{\gamma\delta}
         & =0,
        \notag
    \end{align}
    and likewise for $\hat B$. The left endpoint tensor is
    $V_Be_{AB}\hat A+V_Ae_{BA}\hat B$: its input leg is mapped by the left
    inverse to a pair of virtual indices, which are closed by the domain wall
    and the left action tensor, leaving the output leg and one bond of the
    operator string. The right endpoint tensor is
    $W_Be_{BA}\hat A+W_Ae_{AB}\hat B$. Following
    \cite[Section~III.E]{GarreRubioSchuch2024DomainWall}, the string
    operator $O^{[i,j]}$ on the periodic chain acts by the identity outside
    the sites $i,\dots,j$, by the two endpoint tensors on the sites $i$ and
    $j$, and by the tensor of $O_g$ on the sites between them.
\end{definition}

\begin{lemma}[Domain-wall strings create pairs of domain walls]
    \label{thm:mpug_domain_wall_string_creation}
    \lean{MPOTensor.stringOperator_mulVec_wallConfig,
        MPOTensor.stringOperator_mulVec_trace,
        MPOTensor.mpv_wallConfig,
        MPOTensor.leftAct_mul_physAct_mul_rightAct,
        MPOTensor.leftAct_leftEndpoint,
        MPOTensor.leftAct_leftEndpoint_eq_zero,
        MPOTensor.rightAct_rightEndpoint,
        MPOTensor.rightAct_rightEndpoint_eq_zero,
        MPOTensor.GroupFamily.BlockActionData.wallString_mulVec_mpv_left,
        MPOTensor.GroupFamily.BlockActionData.wallString_mulVec_mpv_right}
    \leanok
    \uses{def:mpug_domain_wall_string, thm:mpug_domain_wall_state}
    On the periodic chain of every length $L\geq2$ and for all sites $i<j$,
    \begin{align}
        O^{[i,j]}\ket{\psi_A}
         & =\ket{\psi(A\text{-}B\text{-}A)},
         &
        O^{[i,j]}\ket{\psi_B}
         & =\ket{\psi(B\text{-}A\text{-}B)},
        \notag
    \end{align}
    with, as in \cite[Section~III.E]{GarreRubioSchuch2024DomainWall}, the
    domain walls $e_{AB}$ at $i$ and $e_{BA}$ at $j$ in the first state and
    $e_{BA}$ at $i$ and $e_{AB}$ at $j$ in the second.
\end{lemma}

\begin{proof}\leanok
    On a state given by the trace of a chain, the string operator acts by
    its endpoint tensors on the letters of the two endpoint sites and by the
    physical action of $O_g$ on the chain between them. On $A$ the left
    inverse $\hat A$ turns the left endpoint into $e_{AB}V_B$ and the right
    endpoint into $W_Be_{BA}$, while the terms with $\hat B$ vanish; the
    action tensors reduce the string of $O_g\cdot A$ between them to $B$.
    On $B$ the terms with $\hat B$ act in the same way.
\end{proof}

\begin{lemma}[Exchange of two domain-wall strings]
    \label{thm:mpug_domain_wall_string_exchange}
    \lean{MPOTensor.wallConfig_resplit,
        MPOTensor.twoWallChain_append,
        MPOTensor.wallConfig_eq_wallConfig,
        MPOTensor.twoWallMPV_wallConfig_comp_cast,
        MPOTensor.fourWallCoeff,
        MPOTensor.exists_wallConfig_wallSplit,
        MPOTensor.GroupFamily.BlockActionData.wallString_outer_mul_inner_mulVec_mpv_apply,
        MPOTensor.GroupFamily.BlockActionData.wallString_inner_mul_outer_mulVec_mpv_apply,
        MPOTensor.GroupFamily.BlockActionData.wallString_mul_wallString_mulVec_mpv}
    \leanok
    \uses{thm:mpug_domain_wall_string_creation, thm:mpug_domain_wall_pair}
    Suppose $g$ carries $e_{AB}$ to $e_{BA}$ with phase $c_{AB}$ and $e_{BA}$
    to $e_{AB}$ with phase $c_{BA}$. For sites $i_2<i_1<j_1<j_2$ whose three
    inner regions are longer than a fixed buffer,
    \begin{align}
        O^{[i_2,j_2]}O^{[i_1,j_1]}\ket{\psi_A}
         & =c_{AB}c_{BA}\,O^{[i_1,j_1]}O^{[i_2,j_2]}\ket{\psi_A}.
        \notag
    \end{align}
    The phase stands on the side where it is acquired; the relation printed
    in \cite[Section~III.F]{GarreRubioSchuch2024DomainWall} has it on the
    other side. The buffer is the blocked form of the local action
    \cite{gap:gs24_domain_wall_nondegenerate}.
\end{lemma}

\begin{proof}\leanok
    Both sides are multiples of the state with the domain walls $e_{AB}$,
    $e_{BA}$, $e_{AB}$, $e_{BA}$ at $i_2,i_1,j_1,j_2$. If the inner string
    acts first, the outer string passes over its two domain walls, and
    Lemma~\ref{thm:mpug_domain_wall_pair} gives the phase $c_{AB}c_{BA}$.
    If the outer string acts first, the inner string acts on the $B$ region
    it created, and Lemma~\ref{thm:mpug_domain_wall_string_creation} gives
    no phase.
\end{proof}

\begin{lemma}[The square of the exchange phase]
    \label{lem:mpug_domain_wall_exchange_phase_square}
    \lean{MPOTensor.exists_twoWallMPV_ne_zero,
        MPOTensor.GroupFamily.BlockActionData.IsDomainWallAction.mul_sq_eq_one_of_mpo_mul_self_eq_one}
    \leanok
    \uses{thm:mpug_domain_wall_state}
    Let $g$ carry $e_{AB}$ to $e_{BA}$ with phase $c_{AB}$ and $e_{BA}$ to
    $e_{AB}$ with phase $c_{BA}$, where $A$ and $B$ are normal, and suppose
    that $U=O_g$ satisfies $U^2=\mathbf 1$ on the periodic chain of every
    length. Then, as in \cite[Section~III.B]{GarreRubioSchuch2024DomainWall},
    $(c_{AB}c_{BA})^2=1$.
\end{lemma}

\begin{proof}\leanok
    By Lemma~\ref{thm:mpug_domain_wall_state},
    $U\ket{\psi(A\text{-}B\text{-}A)}=c_{AB}c_{BA}\ket{\psi(B\text{-}A\text{-}B)}$
    and $U\ket{\psi(B\text{-}A\text{-}B)}=c_{BA}c_{AB}\ket{\psi(A\text{-}B\text{-}A)}$
    whenever the regions are long, so $U^2=\mathbf 1$ gives
    $\ket{\psi(A\text{-}B\text{-}A)}=(c_{AB}c_{BA})^2\ket{\psi(A\text{-}B\text{-}A)}$.
    The state is nonzero for suitable long regions: the words of a fixed
    long length span the full matrix algebras of $A$ and $B$, so if every
    coefficient vanished, then $\operatorname{tr}(Xe_{AB}^iYe_{BA}^j)=0$ for
    all matrices $X$ and $Y$, and choosing matrix units for $X$ and $Y$
    gives $e_{AB}^i=0$ or $e_{BA}^j=0$ for every pair of letters.
\end{proof}

\begin{corollary}[Exchange of two domain-wall strings, printed orientation]
    \label{cor:mpug_domain_wall_string_exchange_printed}
    \lean{MPOTensor.isInjective_of_physPairing_eq_one,
        MPOTensor.GroupFamily.BlockActionData.wallString_mul_wallString_mulVec_mpv_of_mpo_mul_self_eq_one}
    \leanok
    \uses{thm:mpug_domain_wall_string_exchange,
        lem:mpug_domain_wall_exchange_phase_square}
    In the setting of Lemma~\ref{thm:mpug_domain_wall_string_exchange},
    suppose that $U=O_g$ satisfies $U^2=\mathbf 1$ on the periodic chain of
    every length. Then the relation of
    \cite[Section~III.F]{GarreRubioSchuch2024DomainWall} holds:
    \begin{align}
        O^{[i_1,j_1]}O^{[i_2,j_2]}\ket{\psi_A}
         & =c_{AB}c_{BA}\,O^{[i_2,j_2]}O^{[i_1,j_1]}\ket{\psi_A}.
        \notag
    \end{align}
\end{corollary}

\begin{proof}\leanok
    The tensors $A$ and $B$ have left inverses, so their matrices span the
    full matrix algebras and they are normal. Lemma~\ref{lem:mpug_domain_wall_exchange_phase_square}
    gives $(c_{AB}c_{BA})^2=1$. Multiplying the identity of
    Lemma~\ref{thm:mpug_domain_wall_string_exchange} by $c_{AB}c_{BA}$
    gives the result.
\end{proof}

\begin{lemma}[Exchange of two left endpoints of domain-wall strings]
    \label{thm:mpug_domain_wall_endpoint_exchange}
    \lean{MPOTensor.GroupFamily.BlockActionData.wallString_cross_over_right_mulVec_mpv_apply,
        MPOTensor.GroupFamily.BlockActionData.wallString_cross_over_left_mulVec_mpv_apply,
        MPOTensor.GroupFamily.BlockActionData.wallString_mul_wallString_mulVec_mpv_swap_left_far,
        MPOTensor.GroupFamily.BlockActionData.wallString_mul_wallString_mulVec_mpv_swap_left_near}
    \leanok
    \uses{thm:mpug_domain_wall_string_exchange}
    In the setting of Lemma~\ref{thm:mpug_domain_wall_string_exchange},
    let $j<i$ and let $r_1\neq r_2$ be two sites to the right of $i$. If the
    regions between the four sites $j,i,r_1,r_2$ are longer than a fixed
    buffer, then
    \begin{align}
        O^{[j,r_1]}O^{[i,r_2]}\ket{\psi_A}
         & =c_{AB}\,O^{[i,r_1]}O^{[j,r_2]}\ket{\psi_A}.
        \notag
    \end{align}
    This is the exchange of two left endpoints
    $O^{[j]}_xO^{[i]}_y\ket{\psi_A}=c_{AB}\,O^{[i]}_xO^{[j]}_y\ket{\psi_A}$
    of \cite[Section~III.F]{GarreRubioSchuch2024DomainWall}, with the
    half-chain operators glued with the right endpoints of their strings at
    $r_1$ and $r_2$; the identity of half-chain objects with open legs is not
    formalized \cite{gap:gs24_domain_wall_half_string_gluing}.
\end{lemma}

\begin{proof}\leanok
    Every side is a multiple of the state with the domain walls $e_{AB}$,
    $e_{BA}$, $e_{AB}$, $e_{BA}$ at the four sites in increasing order. For
    $r_2<r_1$ the left side is a nested pair whose outer string acts second,
    with phase $c_{AB}c_{BA}$ by
    Lemma~\ref{thm:mpug_domain_wall_string_exchange}; on the right side the
    string on $[i,r_1]$ starts in the $B$ region, passes over the domain wall
    $e_{BA}$ at $r_2$ and ends in the $A$ region, so the local action gives
    $c_{BA}$. For $r_1<r_2$ the string on $[j,r_1]$ on the left side passes
    over the domain wall $e_{AB}$ at $i$, with phase $c_{AB}$, and the right
    side is a nested pair whose inner string acts second, with no phase.
\end{proof}

\begin{lemma}[Rescaling domain walls]
    \label{thm:mpug_domain_wall_rescaling}
    \lean{MPOTensor.actRect_smul,
        MPOTensor.GroupFamily.BlockActionData.IsDomainWallAction.smul_source,
        MPOTensor.GroupFamily.BlockActionData.IsDomainWallAction.smul_target,
        MPOTensor.GroupFamily.BlockActionData.IsDomainWallAction.exists_eq_of_mul_self}
    \leanok
    \uses{def:mpug_domain_wall_action}
    If $g$ carries $e$ to $e'$ with phase $c$ and $a\neq0$, then it carries $a\,e$
    to $e'$ with phase $ac$, and $e$ to $a\,e'$ with phase $c/a$. In
    particular, as in \cite[Section~III.F]{GarreRubioSchuch2024DomainWall},
    if an involution $g$ exchanges $e_{AB}$ and $e_{BA}$ with phases $c_{AB}$
    and $c_{BA}$, then for every $s$ with $s^2=c_{AB}c_{BA}$ the walls
    $e_{AB}$ and $(c_{AB}/s)\,e_{BA}$ are exchanged with both phases equal
    to $s$.
\end{lemma}

\begin{proof}\leanok
    The rectangular action is linear in the domain wall, and rescaling
    $e_{BA}$ by $c_{AB}/s$ gives the last claim.
\end{proof}

\begin{definition}[A family of domain walls permuted by the group]
    \label{def:mpug_domain_wall_family}
    \lean{MPOTensor.GroupFamily.BlockActionData.IsDomainWallFamily,
        MPOTensor.GroupFamily.BlockActionData.IsDomainWallAction.congr_target}
    \leanok
    \uses{def:mpug_domain_wall_action}
    Following \cite[Section~IV.B]{GarreRubioSchuch2024DomainWall}, with
    nonzero walls and phases by convention
    \cite{gap:gs24_domain_wall_nondegenerate}, a family of domain walls
    permuted by the group is a nonzero domain wall $e_{yz}$ between $y$ and
    $z$ for every ordered pair of blocks, together with nonzero phases
    $B^g_{y,z}$ such that every $g$ carries $e_{yz}$ to $e_{gy,gz}$ with
    phase $B^g_{y,z}$.
\end{definition}

\begin{lemma}[The interchange phase of a family of domain walls]
    \label{thm:mpug_domain_wall_interchange}
    \lean{MPOTensor.GroupFamily.BlockActionData.IsDomainWallFamily.mul_eq,
        MPOTensor.GroupFamily.BlockActionData.IsDomainWallFamily.prod_eq_prod_lSymbol,
        MPOTensor.GroupFamily.BlockActionData.IsDomainWallFamily.prod_eq_inv_cyclicInvariant,
        MPOTensor.GroupFamily.BlockActionData.IsDomainWallFamily.prod_eq_domainWallPhase,
        TNLean.Algebra.LSymbol.prod_div_eq_inv_cyclicInvariant,
        TNLean.Algebra.LSymbol.prod_Ico_div_eq_inv_cyclicInvariant}
    \leanok
    \uses{def:mpug_domain_wall_family, def:mpug_permuted_block_l_symbol,
        def:mpug_anomaly_three_cochain}
    As in \cite[Section~IV.C]{GarreRubioSchuch2024DomainWall}, for a family
    of domain walls permuted by the group, every $g$ with $g^n=e$, and every
    block $x$, the interchange phase depends only on the anomaly:
    \begin{align}
        \prod_{k=0}^{n-1}B^g_{g^k(gx),\,g^kx}
         & =\prod_{k=0}^{n-1}\frac{L^{gx}_{g,g^k}}{L^x_{g,g^k}}
        =\Bigl(\prod_{k=0}^{n-1}\omega(g,g^k,g)\Bigr)^{-1}.
        \notag
    \end{align}
    The middle product includes the factor $k=0$, $L^{gx}_{g,e}/L^x_{g,e}$,
    which is one for normalized L-symbols; the source's middle product runs
    over $1\leq k<n$, and for normalized L-symbols that product is again
    $\bigl(\prod_{k=0}^{n-1}\omega(g,g^k,g)\bigr)^{-1}$.
\end{lemma}

\begin{proof}\leanok
    \uses{thm:mpug_domain_wall_pentagon, thm:mpug_permuted_block_l_compatibility}
    With $\beta_k=B^{g^k}_{gx,x}$,
    Lemma~\ref{thm:mpug_domain_wall_pentagon} for $(g,g^k)$ gives
    $B^g_{g^k(gx),g^kx}\,\beta_k=(L^{gx}_{g,g^k}/L^x_{g,g^k})\,\beta_{k+1}$.
    Multiplying over $k<n$ and using $\beta_n=\beta_0$ gives the first
    equality. Compatibility at $(g,g^k,g)$ gives
    $L^{gx}_{g,g^k}/L^x_{g,g^k}=\omega(g,g^k,g)^{-1}
        (L^x_{g,g^{k+1}}/L^x_{g,g^k})(L^x_{g^{k+1},g}/L^x_{g^k,g})^{-1}$,
    and both quotients telescope to one over $k<n$.
\end{proof}

\begin{lemma}[The stabilizer of two blocks acts projectively]
    \label{thm:mpug_domain_wall_projective}
    \lean{TNLean.Algebra.LSymbol.ratioCocycle,
        TNLean.Algebra.LSymbol.ratioCocycle_isMulCocycle₂,
        MPOTensor.GroupFamily.BlockActionData.IsDomainWallFamily.mul_eq_of_fixed,
        MPOTensor.GroupFamily.BlockActionData.IsDomainWallFamily.mul_eq_of_fixed_self}
    \leanok
    \uses{def:mpug_domain_wall_family, def:mpug_permuted_block_l_symbol}
    Let $H_{y,z}$ be the subgroup of elements fixing the blocks $y$ and $z$.
    For $a,b\in H_{y,z}$,
    \begin{align}
        B^a_{y,z}B^b_{y,z}
         & =\alpha(a,b)\,B^{ab}_{y,z},
         &
        \alpha(a,b)
         & =\frac{L^y_{a,b}}{L^z_{a,b}},
        \notag
    \end{align}
    and $\alpha$ is a two-cocycle on $H_{y,z}$; for $y=z$ the phases form a
    linear representation. The source states this for the subgroup fixing
    every block \cite[Section~IV.C]{GarreRubioSchuch2024DomainWall}, which is
    contained in $H_{y,z}$; see \cite{gap:gs24_unbroken_subgroup_stabilizer}.
\end{lemma}

\begin{proof}\leanok
    \uses{thm:mpug_domain_wall_pentagon, thm:mpug_permuted_block_l_compatibility}
    The first identity is Lemma~\ref{thm:mpug_domain_wall_pentagon} with
    $by=y$ and $bz=z$. Dividing the compatibility relations at $y$ and at $z$
    for $(a,b,c)$ with $c\in H_{y,z}$ cancels $\omega(a,b,c)$ and gives the
    two-cocycle equation.
\end{proof}

\begin{definition}[Domain-wall strings of a group element]
    \label{def:mpug_group_domain_wall_string}
    \lean{MPOTensor.GroupFamily.IsSeparatingLeftInverseFamily,
        MPOTensor.GroupFamily.BlockActionData.groupWallLeftEndpoint,
        MPOTensor.GroupFamily.BlockActionData.groupWallRightEndpoint,
        MPOTensor.GroupFamily.BlockActionData.groupWallString,
        MPOTensor.nestedWallCoeff}
    \leanok
    \uses{def:mpug_domain_wall_string, def:mpug_domain_wall_family}
    Let the blocks $A_x$ form a finite family with left inverses $\hat A_x$
    of separated supports: $\hat A_x$ is a left inverse of $A_x$ and
    $\sum_\sigma(\hat A_x^\sigma)_{\alpha\beta}(A_y^\sigma)_{\gamma\delta}=0$
    for $y\neq x$. Let $e_{yz}$ be a family of domain walls permuted by the
    group. Following \cite[Section~IV.B]{GarreRubioSchuch2024DomainWall},
    the left endpoint tensor of $g$ is
    $\sum_xV_{g,x}e_{x,gx}\hat A_x$, the right endpoint tensor is
    $\sum_xW_{g,x}e_{gx,x}\hat A_x$, and the string $O_g^{[i,j]}$ acts by the
    identity outside the sites $i,\dots,j$, by the two endpoint tensors on
    the sites $i$ and $j$, and by the tensor of $O_g$ between them.
\end{definition}

\begin{lemma}[Strings of a group element create pairs of domain walls]
    \label{lem:mpug_group_domain_wall_string_creation}
    \lean{MPOTensor.leftAct_zero, MPOTensor.rightAct_zero,
        MPOTensor.leftAct_sum, MPOTensor.rightAct_sum,
        MPOTensor.GroupFamily.BlockActionData.leftAct_groupWallLeftEndpoint,
        MPOTensor.GroupFamily.BlockActionData.rightAct_groupWallRightEndpoint,
        MPOTensor.GroupFamily.BlockActionData.groupWallString_mulVec_mpv}
    \leanok
    \uses{def:mpug_group_domain_wall_string,
        thm:mpug_domain_wall_string_creation}
    On the periodic chain of every length $L\geq2$ and for all sites $i<j$,
    $O_g^{[i,j]}\ket{\psi_{A_x}}=\ket{\psi(x,gx)}$, the state with the domain
    wall $e_{x,gx}$ at $i$, the block $A_{gx}$ between $i$ and $j$, and the
    domain wall $e_{gx,x}$ at $j$
    \cite[Section~IV.B]{GarreRubioSchuch2024DomainWall}.
\end{lemma}

\begin{proof}\leanok
    On $A_x$ only the term of $\hat A_x$ in each endpoint survives, leaving
    $e_{x,gx}V_{g,x}$ and $W_{g,x}e_{gx,x}$; the action tensors reduce the
    string of $O_g\cdot A_x$ between them to $A_{gx}$.
\end{proof}

\begin{lemma}[Exchange of two domain walls of a group element]
    \label{lem:mpug_group_domain_wall_exchange}
    \lean{MPOTensor.GroupFamily.BlockActionData.groupWallString_outer_mul_inner_mulVec_mpv_apply,
        MPOTensor.GroupFamily.BlockActionData.groupWallString_inner_mul_outer_mulVec_mpv_apply,
        MPOTensor.GroupFamily.BlockActionData.groupWallString_mul_groupWallString_mulVec_mpv,
        MPOTensor.GroupFamily.BlockActionData.groupWallString_cross_mulVec_mpv_apply,
        MPOTensor.GroupFamily.BlockActionData.groupWallString_mul_groupWallString_mulVec_mpv_swap_left}
    \leanok
    \uses{lem:mpug_group_domain_wall_string_creation, thm:mpug_domain_wall_pair}
    In the setting of Definition~\ref{def:mpug_group_domain_wall_string},
    let the regions between the sites considered below be longer than a
    fixed buffer. For sites $i_2<i_1<j_1<j_2$,
    \begin{align}
        O_g^{[i_2,j_2]}O_g^{[i_1,j_1]}\ket{\psi_{A_x}}
         & =B^g_{x,gx}B^g_{gx,x}\,O_g^{[i_1,j_1]}O_g^{[i_2,j_2]}\ket{\psi_{A_x}},
        \notag
    \end{align}
    and for sites $j<i<r_2<r_1$,
    \begin{align}
        O_g^{[j,r_1]}O_g^{[i,r_2]}\ket{\psi_{A_x}}
         & =B^g_{x,gx}\,O_g^{[i,r_1]}O_g^{[j,r_2]}\ket{\psi_{A_x}}.
        \notag
    \end{align}
    The second identity is the exchange of two left endpoints
    $O^{[j]}_{g,a}O^{[i]}_{g,b}\ket{\psi_x}=B^g_{x,gx}\,O^{[i]}_{g,a}O^{[j]}_{g,b}\ket{\psi_x}$
    of \cite[Section~IV.B]{GarreRubioSchuch2024DomainWall}, with the
    half-chain operators glued with the right endpoints of their strings
    \cite{gap:gs24_domain_wall_half_string_gluing}; the buffer is the blocked
    form of the local action \cite{gap:gs24_domain_wall_nondegenerate}.
\end{lemma}

\begin{proof}\leanok
    Every side is a multiple of the state with the domain walls $e_{x,gx}$,
    $e_{gx,g^2x}$, $e_{g^2x,gx}$, $e_{gx,x}$ at the four sites in increasing
    order, between the regions $A_x$, $A_{gx}$, $A_{g^2x}$, $A_{gx}$, $A_x$.
    In a nested pair whose outer string acts second, the outer string passes
    over the domain walls $e_{x,gx}$ and $e_{gx,x}$ created by the inner one,
    with phase $B^g_{x,gx}B^g_{gx,x}$ by
    Lemma~\ref{thm:mpug_domain_wall_pair}. If the outer string acts first,
    the inner string acts on the region $A_{gx}$ and creates the walls
    $e_{gx,g^2x}$ and $e_{g^2x,gx}$ with no phase. In the crossed pair, the
    string on $[i,r_1]$ acting second starts in the region $A_{gx}$, passes
    over the domain wall $e_{gx,x}$ at $r_2$ and ends in the region $A_x$,
    with phase $B^g_{gx,x}$; dividing the two phases gives $B^g_{x,gx}$.
\end{proof}

\begin{lemma}[Domain walls of an involution detect the anomaly]
    \label{thm:mpug_domain_wall_involution}
    \lean{MPOTensor.GroupFamily.BlockActionData.IsDomainWallAction.mul_eq_lSymbol_of_mul_self_eq_one,
        MPOTensor.GroupFamily.BlockActionData.IsDomainWallAction.mul_eq_omega_of_mul_self_eq_one,
        MPOTensor.GroupFamily.BlockActionData.IsDomainWallAction.mul_eq_mul_of_mul_self_eq_one,
        MPOTensor.GroupFamily.BlockActionData.IsDomainWallAction.mpo_mulVec_twoWallMPV_eq_omega,
        MPOTensor.GroupFamily.BlockActionData.IsDomainWallAction.congr_elem}
    \leanok
    \uses{def:mpug_domain_wall_action, def:mpug_permuted_block_l_symbol,
        def:mpug_anomaly_three_cochain}
    Let $g^2=e$ exchange the blocks $x$ and $y$, let $g$ carry the domain wall
    $e_{AB}$ between $x$ and $y$ to $e_{BA}$ with phase $c_{AB}$ and
    $e_{BA}$ to $e_{AB}$ with phase $c_{BA}$. Then
    \begin{align}
        c_{AB}\,c_{BA}
         & =\frac{L^x_{g,g}L^x_{e,e}}{L^y_{g,g}L^y_{e,e}}
        =\omega(g,g,g)\,\omega(g,e,g),
        \notag
    \end{align}
    so $c_{AB}c_{BA}$ does not depend on the action tensors, the fusion
    tensors or the domain walls, and
    $U\ket{\psi(A\text{-}B\text{-}A)}=\omega(g,g,g)\,\omega(g,e,g)
        \ket{\psi(B\text{-}A\text{-}B)}$. When the tensors of the identity are
    trivial, as in \cite[Section~III.C]{GarreRubioSchuch2024DomainWall},
    $L_{e,e}=1$ and $\omega(g,e,g)=1$, which yields
    $c_{AB}c_{BA}=L_A/L_B=\omega$.
\end{lemma}

\begin{proof}\leanok
    \uses{thm:mpug_domain_wall_pentagon, thm:mpug_domain_wall_state,
        thm:mpug_permuted_block_l_compatibility,
        thm:mpug_involution_anomaly_identity_ratio}
    Applying $g$ twice, Lemma~\ref{thm:mpug_domain_wall_pentagon} shows that
    $g^2=e$ carries $e_{AB}$ to itself with phase
    $b=(L^y_{g,g}/L^x_{g,g})\,c_{BA}c_{AB}\neq0$. The same theorem for $(e,e)$
    gives $b^2=(L^x_{e,e}/L^y_{e,e})\,b$, so $b=L^x_{e,e}/L^y_{e,e}$, which is
    the first equality. The second equality is
    Lemma~\ref{thm:mpug_involution_anomaly_identity_ratio} for the compatible
    L-symbols of Lemma~\ref{thm:mpug_permuted_block_l_compatibility}. Its
    right-hand side involves only the fusion tensors, and it is invariant
    under their gauge action, so the product is independent of all choices.
    Finally, by Lemma~\ref{thm:mpug_domain_wall_state},
    $U\ket{\psi(A\text{-}B\text{-}A)}=c_{AB}c_{BA}\ket{\psi(B\text{-}A\text{-}B)}
        =\omega(g,g,g)\,\omega(g,e,g)\ket{\psi(B\text{-}A\text{-}B)}$.
\end{proof}

\subsection{Detecting the anomaly with truncated symmetries}

\begin{lemma}[The commutator of two unitaries on a vector]
    \label{lem:mpug_truncation_commutator}
    \lean{Matrix.groupCommutator,
        Matrix.groupCommutator_mulVec_eq_smul,
        Matrix.star_dotProduct_groupCommutator_mulVec}
    \leanok
    Let $U_1$ and $U_2$ be unitary, let $\psi$ be a vector and $c$ a scalar
    with $U_2U_1\psi=c\,U_1U_2\psi$. Then
    \begin{align}
        U_2^\dagger U_1^\dagger U_2U_1\psi
         & =c\,\psi,
         &
        \langle\psi|\,U_2^\dagger U_1^\dagger U_2U_1\,|\psi\rangle
         & =c\,\langle\psi|\psi\rangle.
        \notag
    \end{align}
    For two truncations $U_1=U^{[i_1,j_1]}$ and $U_2=U^{[i_2,j_2]}$ of a
    symmetry whose exchange relation on a ground state has the phase
    $c=\omega$, this is the last step of the detection of the anomaly in
    \cite[Section~V]{GarreRubioSchuch2024DomainWall}: the expectation value of
    the commutator measures the exchange phase.
\end{lemma}

\begin{proof}\leanok
    Substituting the exchange relation and canceling $U_1^\dagger U_1$ and
    $U_2^\dagger U_2$ gives both identities.
\end{proof}

\begin{lemma}[Truncations of an on-site symmetry]
    \label{thm:mpug_onsite_truncation_commutator}
    \lean{Matrix.onSiteTruncation,
        Matrix.onSiteTruncation_univ,
        Matrix.onSiteTruncation_conjTranspose_mul_self,
        Matrix.groupCommutator_onSiteTruncation,
        Matrix.star_dotProduct_groupCommutator_onSiteTruncation_mulVec}
    \leanok
    Let $u$ satisfy $u^\dagger u=\mathbf 1$ and, for a set $S$ of sites of a
    chain, let $U^S=\bigotimes_{i\in S}u^{[i]}$ act by $u$ on the sites of $S$
    and by the identity elsewhere. For any two sets $S_1$ and $S_2$,
    \begin{align}
        (U^{S_2})^\dagger(U^{S_1})^\dagger U^{S_2}U^{S_1}
         & =\mathbf 1,
        \notag
    \end{align}
    so, as in \cite[Section~V]{GarreRubioSchuch2024DomainWall} for the
    trivial anomaly of on-site symmetries, the expectation value of the
    commutator on every normalized state is $1$.
\end{lemma}

\begin{proof}\leanok
    The four operators are tensor products over the sites. On a site in both
    sets the product of the four factors is
    $u^\dagger u^\dagger uu=\mathbf 1$, on a site in one set it is
    $u^\dagger u=\mathbf 1$, and elsewhere all four factors are the
    identity.
\end{proof}

\begin{lemma}[The detector depends on the endpoints of the truncations]
    \label{thm:mpug_truncation_detector_endpoint_dependence}
    \lean{Matrix.finKronecker_ite_smul_one,
        Matrix.dressed_onSiteTruncation_conjTranspose_mul_self,
        Matrix.groupCommutator_dressed_onSiteTruncation_pauliX,
        Matrix.star_dotProduct_groupCommutator_dressed_onSiteTruncation_pauliX}
    \leanok
    Let $X$ and $Z$ be the Pauli matrices, and for sets of sites $S_1$, $S_2$
    of a chain and sites $j_1\in S_2$, $j_2\notin S_1$ with $j_1\neq j_2$ put
    $U_1=Z_{j_1}X^{S_1}$ and $U_2=Z_{j_2}X^{S_2}$. Then $U_1$ and $U_2$ are
    unitary and
    \begin{align}
        U_2^\dagger U_1^\dagger U_2U_1
         & =-\mathbf 1,
        \notag
    \end{align}
    so the expectation value of the commutator on every normalized state is
    $-1$, although the on-site symmetry $X^{\otimes N}$ has trivial anomaly.
\end{lemma}

\begin{proof}\leanok
    The four operators are tensor products over the sites. At $j_1$ the
    product of the four factors is $X(XZ)X(ZX)=ZXZX=-\mathbf 1$ if $j_1\in S_1$,
    and $XZXZ=-\mathbf 1$ otherwise. At $j_2$ only $U_2$ acts, and the product is
    $(ZX)^\dagger(ZX)=\mathbf 1$ or $Z^\dagger Z=\mathbf 1$. At every other site
    the factors are $X$ or the identity, and they cancel.
\end{proof}

\begin{remark}[Circuit truncations with a nontrivial detector]
    \label{rem:mpug_truncation_detector_circuit}
    \uses{thm:mpug_truncation_detector_endpoint_dependence}
    The symmetry $X^{\otimes N}$ on a periodic chain of even length is the
    depth-two circuit whose first layer applies $X\otimes ZX$ to the pairs
    $(2k,2k+1)$ and whose second layer applies $Z\otimes\mathbf 1$ to the pairs
    $(2k+1,2k+2)$, since $Z\cdot ZX=X$. For $i$ even and $j$ odd, dropping the
    gates not contained in $[i,j]$ leaves both factors on every odd site of
    $[i,j]$ except $j$, which keeps only $ZX$, so the truncation is
    $Z_jX^{\otimes[i,j]}$. For $i_2<i_1<j_1<j_2$ of these parities,
    Lemma~\ref{thm:mpug_truncation_detector_endpoint_dependence} applies.
    The expectation value of the commutator of two circuit truncations on a
    symmetry-broken ground state is therefore not determined by the anomaly,
    and the detection asserted in
    \cite[Section~V]{GarreRubioSchuch2024DomainWall} fails for an arbitrary
    circuit truncation \cite{gap:gs24_truncation_detector_endpoint_dependence}.
    The identification of the circuit truncation with $Z_jX^{\otimes[i,j]}$
    is not formalized.
\end{remark}

\begin{lemma}[Detection for truncations acting as domain-wall strings]
    \label{thm:mpug_truncation_detector_wall_strings}
    \lean{MPOTensor.GroupFamily.BlockActionData.star_dotProduct_groupCommutator_mulVec_mpv_of_wallString,
        MPOTensor.GroupFamily.BlockActionData.star_dotProduct_groupCommutator_mulVec_mpv_eq_omega}
    \leanok
    \uses{lem:mpug_truncation_commutator, thm:mpug_domain_wall_string_exchange,
        thm:mpug_domain_wall_involution}
    In the setting of Lemma~\ref{thm:mpug_domain_wall_string_exchange}, let
    $O_1=O^{[i_1,j_1]}$ and $O_2=O^{[i_2,j_2]}$ be domain-wall strings on
    sites $i_2<i_1<j_1<j_2$ whose three inner regions are longer than a fixed
    buffer, and let $U_1$ and $U_2$ be unitary operators with
    \begin{align}
        U_2U_1\ket{\psi_A} & =O_2O_1\ket{\psi_A}, &
        U_1U_2\ket{\psi_A} & =O_1O_2\ket{\psi_A}.
        \notag
    \end{align}
    Then
    \begin{align}
        \bra{\psi_A}U_2^\dagger U_1^\dagger U_2U_1\ket{\psi_A}
         & =c_{AB}c_{BA}\braket{\psi_A}{\psi_A}.
        \notag
    \end{align}
    If in addition $g$ is an involution of a normal representation, then
    $c_{AB}c_{BA}=\omega(g,g,g)\,\omega(g,e,g)$. This is the detection of the
    anomaly of \cite[Section~V]{GarreRubioSchuch2024DomainWall} restricted to
    unitaries whose two ordered products on $\ket{\psi_A}$ agree with the
    corresponding products of the domain-wall strings. These identities are
    expected for truncations whose walls on $\ket{\psi_B}$ are the images under
    $U$ of their walls on $\ket{\psi_A}$, the property the source's argument
    uses; that implication is not formalized. By
    Lemma~\ref{thm:mpug_truncation_detector_endpoint_dependence} some
    restriction of this kind is needed
    \cite{gap:gs24_truncation_detector_endpoint_dependence}.
\end{lemma}

\begin{proof}\leanok
    \uses{lem:mpug_truncation_commutator, thm:mpug_domain_wall_string_exchange,
        thm:mpug_domain_wall_involution}
    By Lemma~\ref{thm:mpug_domain_wall_string_exchange} and the hypotheses,
    $U_2U_1\ket{\psi_A}=c_{AB}c_{BA}\,U_1U_2\ket{\psi_A}$, and
    Lemma~\ref{lem:mpug_truncation_commutator} gives the first identity. The
    second is Lemma~\ref{thm:mpug_domain_wall_involution}.
\end{proof}
